% Architecture -> discovery attack -> measured findings -> limitations.
% About 12 minutes plus questions. Explanations and provenance live in \note{}.
% Build from this directory; see README.md. Research notes live in the sibling support folder.
\documentclass[aspectratio=169]{beamer}
\usetheme{Chalmers}
\usepackage[utf8]{inputenc}
\usepackage[T1]{fontenc}
\usepackage[british]{babel}
\usepackage{graphicx}
\usepackage{array}
\usepackage{tikz}
\usepackage{pgfpages}
\usetikzlibrary{positioning,arrows.meta,shapes.geometric,fit,backgrounds,calc}
\title{Who decides who you train with?}
\subtitle{Hivemind control-plane security}
\author{Sarat Yellapragada}
\institute{Chalmers University of Technology}
\date{2026-09-18}
\shorttitle{\ChalmersSectionName}
\chalmersheadline{logo}
\chalmerssectiondividers{none}
\chalmersfooterpages{both}
% This template version compares switch macros without expanding them.
% Apply these two choices explicitly, without changing the shared theme.
\AtBeginSection[]{}
\renewcommand{\Chalmersinsertpagenumber}{\insertframenumber\,/\,\insertmainframenumber}
\setbeamerfont{frametitle}{size=\fontsize{19}{22}\selectfont,series=\bfseries}
\setlength{\Chalmersframetitlebottom}{0mm}
\ifdefined\presenternotes
  \setbeameroption{show notes on second screen=right}
\else
  \setbeameroption{hide notes}
\fi
\colorlet{ctl}{ChalmersRed}
\colorlet{dat}{ChalmersGreen}
\colorlet{ink}{ChalmersBlack}
\colorlet{muted}{ChalmersGrey}
\tikzset{
  >={Stealth[length=2mm]},
  every node/.style={font=\small,text=ink},
  labeltext/.style={font=\small,align=center},
  quiet/.style={font=\footnotesize,text=muted,align=center},
  heading/.style={font=\normalsize\bfseries,align=center},
  box/.style={rounded corners=2mm,draw=ChalmersPurple!25,fill=ChalmersPurple!5,
    align=center,inner sep=3mm},
  peer/.style={circle,draw=ChalmersPurple,fill=ChalmersPurple!10,
    minimum size=9mm,inner sep=0pt,font=\small\bfseries,text=ChalmersPurple},
  victim/.style={peer,fill=white,line width=1.1pt,double,double distance=1pt},
  attacker/.style={diamond,aspect=1,draw=ctl,fill=ctl!12,
    minimum size=10mm,inner sep=0pt,font=\small\bfseries,text=ctl},
  flow/.style={->,line width=1.1pt,draw=ChalmersPurple},
  control/.style={->,line width=1.1pt,draw=ctl},
  dataflow/.style={<->,line width=1.2pt,draw=dat},
}
% A fixed canvas keeps diagrams at a consistent scale, without shrinking text.
\newenvironment{visual}{\centering\begin{tikzpicture}[x=1cm,y=1cm]
  \path[use as bounding box] (0,0) rectangle (14.4,5.6);
}{\end{tikzpicture}}
% Links stay in the lower margin so evidence tables do not crowd the visuals.
\newcommand{\resultlink}[2]{%
  \begin{tikzpicture}[remember picture,overlay]
    \node[anchor=south east,inner sep=0pt] at ([xshift=-8mm,yshift=7mm]current page.south east)
      {\hyperlink{#1}{\beamergotobutton{#2}}};
  \end{tikzpicture}%
}
\begin{document}

% 1 --------------------------------------------------------------------------
\begin{chalmerscover}
\begin{visual}
  \node[anchor=west,text=white,font=\fontsize{27}{31}\selectfont\bfseries,
    align=left] at (0.2,4.3) {Hivemind};
  \node[anchor=west,text=white,font=\large] at (0.2,2.7)
    {};
  \node[anchor=west,text=white,align=left,font=\small] at (0.2,0.95)
    {Sarat Yellapragada\\[2mm]Secura Lab\\[2mm]Chalmers University of Technology\quad\textbar\quad18 September 2026};
  % \draw[white!45,line width=1pt] (10.3,4.7)--(12.8,3.8)--(11.1,2.5)--cycle;
  % \node[circle,draw=white,fill=ChalmersPurple,minimum size=13mm,
  %   text=white,font=\large\bfseries] at (10.3,4.7) {V};
  % \node[diamond,draw=white,fill=white!18!ChalmersPurple,minimum size=15mm,
  %   text=white,font=\large\bfseries] at (12.8,3.8) {A};
  % \node[circle,draw=white,fill=ChalmersPurple,minimum size=13mm] at (11.1,2.5) {};
  % \draw[white,dashed,rounded corners=5mm] (9.3,1.65) rectangle (13.8,5.55);
\end{visual}
\note{Hivemind trains across independent machines. My project asks how an attacker can
influence which machines average together, and what that does to learning. Walk through
architecture first, then the experimental progress and results. V marks the victim;
A marks an attacker.}
\end{chalmerscover}

\section{Architecture}
% Overview and four visual steps for a network-security audience with basic FL/ML knowledge.
% Included by hivemind-control-plane.tex; shares its styles and speaker notes.

\begin{frame}{Inside a Hivemind peer}
\begin{visual}
  % An overview of components, with only the language boundary connected.
  \fill[ChalmersPurple!4,rounded corners=3mm]
    (1.8,2.1) rectangle (12.6,5.25);
  \node[anchor=west,font=\small\bfseries,text=ChalmersPurple]
    at (2.2,4.94) {Hivemind\quad\textbar\quad Python};

  \foreach \xx/\yy/\name/\purpose in {
    4.55/4.23/{All-reduce}/{Averages group gradients.},
    9.85/4.23/{Optimizer}/{Schedules averaging.},
    4.55/3.2/{Hivemind DHT}/{Stores peer metadata.},
    9.85/3.2/{Matchmaking}/{Forms temporary groups.}} {
    \node[box,draw=none,fill=white,minimum width=4.6cm,minimum height=0.85cm,
      inner sep=0pt] at (\xx,\yy) {};
    \node[font=\small\bfseries,text=ChalmersPurple] at (\xx,\yy+0.17) {\name};
    \node[quiet,text=ink] at (\xx,\yy-0.18) {\purpose};
  }

  \node[box,draw=none,fill=ChalmersPurple!9,minimum width=7.8cm,
    minimum height=0.42cm,inner sep=0pt,font=\footnotesize,text=ChalmersPurple]
    at (7.2,2.43) {P2P wrapper + bindings};

  \draw[<->,ChalmersPurple!75,line width=1pt] (7.2,2.17)--(7.2,1.37);
  \node[anchor=west,font=\footnotesize,text=muted]
    at (7.4,1.77) {Unix-domain sockets};

  % Nesting denotes a library inside the Go daemon, not another process.
  \fill[ChalmersCyan!7,rounded corners=3mm]
    (3.1,0.17) rectangle (11.3,1.32);
  \node[font=\small\bfseries,text=ChalmersPurple] at (7.2,1.05)
    {p2pd\quad\textbar\quad Go};
  \node[box,draw=none,fill=white,minimum width=7.4cm,minimum height=0.5cm,
    inner sep=0pt] at (7.2,0.57) {};
  \node[font=\small\bfseries,text=ChalmersPurple] at (4.8,0.57) {libp2p};
  \node[quiet,text=ink] at (8.3,0.57) {Connects peers securely.};
\end{visual}
\note{This is a map of components; the following slides explain their interactions.
The optimizer coordinates local training and averaging. The DHT stores discovery and
progress metadata, matchmaking forms temporary groups, and all-reduce computes a
weighted gradient average. The shared networking interface is Python's P2P wrapper
and daemon bindings. Python components can span several processes.
\par p2pd is a separate Go process containing libp2p. The single arrow represents local
RPCs and streams over Unix-domain sockets: Python connects to the daemon, and incoming
streams can reach a Python client listener. Lower-level bindings also support IPv4/TCP;
Hivemind creates Unix sockets by default. libp2p handles authenticated, encrypted peer
connections. Gradients bypass DHT storage.
\par Sources: p2p/p2p\_daemon.py:42,165--166,227,255;\\
p2p/p2p\_daemon\_bindings/control.py:57--64,157--165; setup.py:113--116;
progress\_tracker.py:220,247; optimizer.py:528,613--630;
key\_manager.py:62--80; averager.py:428--464,543--554.}
\end{frame}

\begin{frame}{Hivemind's training loop}
\begin{visual}
  \node[heading] at (2.3,5.15) {1\quad Train locally};
  \node[heading] at (7.3,5.15) {2\quad Average in a group};
  \node[heading] at (12.15,5.15) {3\quad Update locally};
  \draw[ChalmersPurple!45,fill=ChalmersPurple!3,rounded corners=3mm]
    (0.3,1.65) rectangle (4.3,4.55);
  \node[heading,text=ChalmersPurple] at (2.3,4.1) {Peer 1};
  \node[box,draw=muted!55,fill=white,inner sep=1.5mm,minimum width=12mm,
    minimum height=8mm] (data) at (1.25,3.0) {data};
  \node[box,fill=white,inner sep=1.5mm,minimum width=14mm,minimum height=10mm,
    font=\small] (model) at (3.15,3.0) {full\\model};
  \draw[flow] (data)--(model);
  \node[quiet] at (2.3,2.12) {compute a local gradient};

  \draw[ChalmersPurple!50,dashed,rounded corners=4mm]
    (5.15,1.65) rectangle (9.55,4.55);
  \node[peer,minimum size=11mm] (p1) at (6.15,3.8) {P1};
  \node[peer,minimum size=11mm] (p2) at (8.55,3.8) {P2};
  \node[peer,minimum size=11mm] (p3) at (7.35,2.4) {P3};
  \draw[dataflow] (p1)--(p2);
  \draw[dataflow] (p2)--(p3);
  \draw[dataflow] (p3)--(p1);
  \node[quiet,text=dat] at (7.35,1.23) {peers exchange gradients};
  \draw[flow] (4.4,3.1)--(5.05,3.1);

  \draw[ChalmersPurple!45,fill=ChalmersPurple!3,rounded corners=3mm]
    (10.35,1.65) rectangle (14.0,4.55);
  \node[heading,text=ChalmersPurple] at (12.175,4.1) {Peer 1};
  \node[box,draw=dat!65,fill=white,inner sep=2mm,minimum width=22mm,
    minimum height=10mm] (updated) at (12.175,3.0) {updated model};
  \draw[->,dat,line width=1.2pt] (9.7,3.0)--(updated.west);
  \node[quiet] at (12.175,2.12) {apply average gradient};

  \draw[flow,rounded corners=2mm] (12.175,1.55)--(12.175,0.65)--(2.3,0.65)--(2.3,1.55);
  \node[font=\footnotesize,text=ChalmersPurple,fill=white,inner sep=1mm] at (7.2,0.65) {repeat};
  \node[quiet] at (7.2,0.05) {Every peer keeps its own data and full model\quad\textbar\quad no central aggregation server};
\end{visual}
\note{Follow Peer 1 from left to right; it is the same machine in all three stages.
First, its own data and full model produce a local gradient, a vector describing a training
update. Second, it joins ready peers in a temporary group and exchanges gradients with them.
The circles P1,P2,P3 are Peers 1,2,3, all doing the same job. Third, Peer 1 applies the group
average to its local model and continues training. Its groupmates do the same. This is the
basic repeated operation Hivemind enables without a central aggregation server.
\par Unlike a server-based FL diagram, the middle is a group of training peers, not an extra
service. Later averaging operations assemble groups again, although membership need not change.
Different groups need not have identical model states. Raw data stays in the local training
pipeline; this does not guarantee gradient privacy. Bootstrap contacts help enter the network,
not coordinate learning. Sources: HIVEMIND\_EXPLAINED sections 1,4,6--8; HANDOFF 4.10.}
\end{frame}

% \begin{frame}{Hivemind's training loop}
% \begin{visual}
%   \node[heading] at (2.2,5.15) {Train at different speeds};
%   \node[heading] at (6.65,5.15) {One shared target};
%   \node[heading] at (11.55,5.15) {Average + update};

%   % Illustrative local sample contributions, not individual batch quotas.
%   \foreach \name/\yy/\width/\samples in {
%     P1/4.45/2.4/600,P2/3.86/1.6/400,P3/3.27/0.8/200} {
%     \node[font=\footnotesize\bfseries,text=ChalmersPurple,anchor=east]
%       at (0.83,\yy) {\name};
%     \fill[ChalmersCyan!65,rounded corners=0.7mm]
%       (1.0,\yy-0.13) rectangle (1.0+\width,\yy+0.13);
%     \node[font=\footnotesize,text=muted,anchor=west]
%       at (3.58,\yy) {\samples};
%   }
%   \draw[flow] (4.45,3.86)--(5.15,3.86);
%   \node[box,draw=none,fill=ChalmersPurple!6,minimum width=2.65cm,
%     minimum height=1.75cm,inner sep=0pt] at (6.65,3.86) {};
%   \node[font=\LARGE\bfseries,text=ChalmersPurple] at (6.65,4.09) {1,200};
%   \node[quiet] at (6.65,3.53) {samples / epoch};
%   \draw[flow] (8.05,3.86)--(8.8,3.86);

%   \draw[ChalmersPurple!40,dashed,rounded corners=3mm]
%     (8.97,2.99) rectangle (14.0,4.78);
%   \node[peer,minimum size=8mm] (t1) at (9.9,4.24) {P1};
%   \node[peer,minimum size=8mm] (t2) at (13.07,4.24) {P2};
%   \node[peer,minimum size=8mm] (t3) at (11.48,3.52) {P3};
%   \draw[dataflow,line width=0.8pt] (t1)--(t2);
%   \draw[dataflow,line width=0.8pt] (t2)--(t3);
%   \draw[dataflow,line width=0.8pt] (t3)--(t1);
%   \node[quiet] at (2.2,2.78) {local samples accumulated};
%   \node[quiet] at (11.55,2.78) {temporary averaging group};
%   \draw[flow] (12.8,2.39)--(1.6,2.39);
%   \node[font=\footnotesize,text=ChalmersPurple,fill=white,inner sep=1mm]
%     at (7.2,2.39) {Next epoch: train + form groups again};

%   % Optional redistribution is distinct from the normal epoch cycle above.
%   \node[font=\small\bfseries,text=ChalmersPurple] at (1.95,1.5) {Optional sharding};
%   \node[quiet] at (1.95,1.04) {off by default};
%   \draw[ChalmersPurple!40,dashed,rounded corners=2mm]
%     (3.65,0.65) rectangle (6.85,1.88);
%   \node[quiet,font=\scriptsize,text=muted,fill=white,inner sep=0.5mm]
%     at (5.25,1.88) {one group};
%   \foreach \xx/\member in {4.17/1,4.89/2,5.61/3,6.33/4}
%     \node[peer,minimum size=5.5mm,font=\scriptsize\bfseries] at (\xx,1.24) {\member};
%   \draw[flow] (7.0,1.24)--(7.85,1.24);
%   \foreach \xx/\bits/\member in {8.65/00/1,10.15/01/2,11.65/10/3,13.15/11/4} {
%     \node[box,draw=ChalmersPurple!25,fill=ChalmersPurple!3,minimum width=1.12cm,
%       minimum height=1.28cm,inner sep=0pt] at (\xx,1.24) {};
%     \node[font=\footnotesize\ttfamily,text=muted] at (\xx,1.61) {\bits};
%     \node[peer,minimum size=5.5mm,font=\scriptsize\bfseries] at (\xx,1.02) {\member};
%   }
%   \node[quiet] at (10.9,0.25) {different discovery buckets next round};
%   \node[quiet] at (3.7,0.25) {Default: groupmates may repeat};
% \end{visual}
% \note{An epoch is a shared training-sample target, not a data pass or a per-peer quota.
% Faster peers contribute more minibatches: 600+400+200=1200 is an illustration.
% Peers retain local data and full models. The progress tracker exchanges sample counts
% and speeds through the DHT, estimates epoch end, and schedules matchmaking near it;
% this is not an exact barrier. Temporary groups average gradients, peers update, then
% later averaging operations run matchmaking again. Their membership may repeat.
% \par The separate four-peer example shows optional sharding. Group assembly assigns
% members distinct next discovery buckets using the shared group ID and member positions.
% The example has enough suffix bits for four distinct assignments. These boxes are
% rendezvous buckets, not assembled groups. Ordinary co-members separate next round;
% an attacker can compute the same assignments and follow. Empty group bits make this
% a no-op by default. Next keys change on assembly, not at a global epoch barrier.
% \par Sources: optimizer.py:54--67,559--567; progress\_tracker.py:128--149,308--322;
% averager.py:120,427--459,501--512; key\_manager.py:94--104;
% matchmaking.py:387,404; HANDOFF 4.4,4.8.}
% \end{frame}

\begin{frame}{How a group forms}
\begin{visual}
  \node[heading,text=ChalmersPurple] at (1.55,5.2) {Peer 1};
  \node[box,draw=ctl!50,fill=ctl!6,inner sep=2mm,font=\small\bfseries]
    at (5.45,5.2) {DHT replicas};
  \node[heading,text=ChalmersPurple] at (9.1,5.2) {Peer 2};
  \node[heading,text=ChalmersPurple] at (12.85,5.2) {Peer 3};
  \foreach \xx in {1.55,5.45,9.1,12.85}
    \draw[muted!45,dashed] (\xx,4.8)--(\xx,0.65);
  \draw[->,muted!65] (0.3,4.45)--(0.3,1.1);
  \node[quiet,rotate=90] at (0.08,2.75) {time};
  \draw[control] (9.1,4.25)--node[above,font=\footnotesize]{availability + deadline}(5.45,4.25);
  \draw[control] (1.55,3.6)--node[above,font=\footnotesize]{lookup job / bucket}(5.45,3.6);
  \draw[control] (5.45,2.95)--node[above,font=\footnotesize]{candidates + deadlines}(1.55,2.95);
  \node[quiet,text=ctl,align=left,anchor=west] at (9.3,3.3) {earliest visible deadline\\becomes candidate leader};
  \filldraw[fill=ctl!12,draw=ctl!60] (8.96,0.7) rectangle (9.24,2.6);
  \draw[control] (1.55,2.25)--node[above,font=\footnotesize]{join request}(8.96,2.25);
  \draw[control] (12.85,1.75)--node[above,font=\footnotesize]{join request}(9.24,1.75);
  \draw[control] (8.96,1.2)--node[above,font=\footnotesize]{member list + group ID}(1.55,1.2);
  \draw[control] (9.24,0.85)--node[above,font=\footnotesize]{member list}(12.85,0.85);
  \node[box,draw=none,fill=ChalmersCream,minimum width=13.6cm,inner sep=2mm] at (7.2,0.15)
    {The temporary leader checks requests and admits members.};
\end{visual}
\note{Read this as a sequence diagram, downward in time. Only selected messages are shown;
all searching peers can publish advertisements. A job and optional bucket identify a DHT
rendezvous key. Records advertise a PeerID, availability and an expiry that also serves as
the intended matchmaking deadline. Peers look up candidates and prefer the earliest visible
deadline, with PeerID as tie-breaker. There is no consensus election or guaranteed common view.
\par Peer 2 is the candidate leader in this example. Others send it direct join RPCs. It checks
compatibility and protocol fields, admits followers, then distributes the group ID and member
list when full or when its deadline permits. The leader assembles the group; it is not the
sole gradient reducer. Groups form again for later averaging. Moving to different DHT buckets
is optional: the default empty group bits leave the rendezvous key unchanged. Sources:
key\_manager.py:43--104; matchmaking.py:262--405,469--498; HANDOFF sections 1,4.3,4.8.}
\end{frame}

\begin{frame}{All-reduce: averaging within the group}
\begin{visual}
  \node[heading] at (2.05,5.15) {1\quad Split};
  \node[heading] at (7.2,5.15) {2\quad Average each slice};
  \node[heading] at (12.35,5.15) {3\quad Reassemble};
  \foreach \peer/\yy/\a/\b/\c in {1/4.05/2/4/6,2/2.95/4/6/8,3/1.85/6/8/10} {
    \node[quiet,anchor=east] at (0.95,\yy) {Peer \peer};
    \foreach \xx/\val/\col in {1.45/\a/ChalmersCyan,2.25/\b/ChalmersGreen,3.05/\c/ChalmersPurple} {
      \node[draw=\col!65,fill=\col!12,minimum width=7.5mm,minimum height=7mm,
        inner sep=0pt,font=\normalsize\bfseries] at (\xx,\yy) {\val};
    }
    \node[quiet,anchor=east] at (11.3,\yy) {Peer \peer};
    \foreach \xx/\val/\col in {11.8/4/ChalmersCyan,12.6/6/ChalmersGreen,13.4/8/ChalmersPurple} {
      \node[draw=\col!65,fill=\col!12,minimum width=7.5mm,minimum height=7mm,
        inner sep=0pt,font=\normalsize\bfseries] at (\xx,\yy) {\val};
    }
  }
  \foreach \peer/\yy/\vals/\mean/\col in {
    1/4.05/{2,4,6}/4/ChalmersCyan,
    2/2.95/{4,6,8}/6/ChalmersGreen,
    3/1.85/{6,8,10}/8/ChalmersPurple} {
    \node[box,draw=\col!65,fill=\col!8,minimum width=4.0cm,minimum height=8mm,
      inner sep=2mm] at (7.2,\yy)
      {\textbf{Peer \peer}\quad $\operatorname{mean}(\vals)=\mean$};
  }
  \draw[dataflow,->] (3.65,2.95)--(4.85,2.95);
  \draw[dataflow,->] (9.45,2.95)--(10.1,2.95);
  \node[quiet] at (2.05,0.95) {local gradients};
  \node[quiet] at (7.2,0.95) {a different reducer for each slice};
  \node[quiet] at (12.45,0.95) {same group average};
  \node[quiet] at (7.2,0.15)
    {All-reduce: every group member obtains the aggregate\quad\textbar\quad equal-weight illustration};
\end{visual}
\note{All-reduce means all group members obtain the aggregate. The toy vectors make the
operation concrete: each peer contributes three numbers, and the group average is [4,6,8].
Peer 1 reduces the blue slice, Peer 2 the green slice, Peer 3 the purple slice. Each sender
contacts the assigned reducer directly; this is not ring forwarding. The same group carries
out all three stages. The DHT stores metadata; peers exchange gradient slices directly using
libp2p, which may use transport relays underneath.
\par Real gradients are long vectors, slices can have unequal sizes, and Hivemind uses a
weighted mean. Reducers return averaged-minus-local deltas; peers reconstruct the result and
the optimizer applies the update. Equal slices and equal weights here are explanatory choices.
The temporary group leader is not a central aggregation server. Sources: allreduce.py:93--111,
164--216,335--356; partition.py:242--260; HANDOFF 4.10--4.11.}
\end{frame}

\begin{frame}{Two IDs}
\begin{visual}
  \node[heading] at (3.5,5.12) {Place records};
  \node[heading] at (10.9,5.12) {Contact peers};
  \node[box,draw=ctl!40,fill=ctl!5,minimum width=5.8cm,minimum height=2.75cm]
    at (3.5,3.3) {};
  \node[box,minimum width=5.8cm,minimum height=2.75cm] at (10.9,3.3) {};
  \node[heading,text=ctl] at (3.5,4.28) {Hivemind DHT: node\_id};
  \node[heading,text=ChalmersPurple] at (10.9,4.28) {libp2p: PeerID};

  \node[labeltext] at (3.5,3.5) {160-bit position in the DHT};
  \node[quiet] at (3.5,2.85) {By default, the 5 nearest nodes\\(XOR distance) store each record.};
  \node[labeltext,text=ctl,font=\small\bfseries] at (3.5,2.25) {chosen by the peer};
  \node[labeltext] at (10.9,3.5) {identity derived from a public key};
  \node[quiet] at (10.9,2.85) {libp2p resolves its address\\and authenticates the connection};
  \node[labeltext,text=ChalmersPurple,font=\small\bfseries] at (10.9,2.25) {key ownership checked};

  % Both identifiers belong to the same peer; there is no ID-conversion arrow.
  \draw[muted!35,rounded corners=2mm,fill=white] (4.0,0.43) rectangle (10.05,1.78);
  \node[font=\footnotesize\bfseries,text=muted] at (7.025,1.5) {One peer, two IDs};
  \node[box,draw=muted!40,fill=white,inner sep=2mm,minimum width=2.2cm]
    (key) at (1.6,0.88) {record key};
  \node[box,draw=ctl!55,fill=ctl!5,inner sep=2mm,minimum width=2.9cm]
    (pos) at (5.62,0.88) {nearest \texttt{node\_id}};
  \node[box,inner sep=2mm,minimum width=2.1cm]
    (endpoint) at (8.75,0.88) {PeerID};
  \node[box,inner sep=2mm,minimum width=2.5cm]
    (rpc) at (12.5,0.88) {connect + RPC};
  \draw[control] (key)--(pos);
  \draw[flow] (endpoint)--(rpc);
  \node[quiet] at (7.2,0.02)
    {The routing table stores both IDs; libp2p connects using PeerID.};
\end{visual}
\note{The DHT uses node\_id as a 160-bit XOR coordinate for routing and replica placement.
It normally stores each record on the five nearest nodes, trying other candidates if stores
fail. Additional cached copies can exist.
\par PeerID identifies a libp2p endpoint through its public key. libp2p resolves network
addresses and authenticates key ownership; PeerID itself is not an IP address.
Hivemind builds its placement overlay on this networking layer, so the IDs have different jobs.
The routing table stores both for each peer: select a node\_id, then use its associated
PeerID to connect. One ID is not computed from the other. The diagram shows one holder;
real lookups and replication involve several nodes.
\par Separate roles do not require unrelated IDs. Here node\_id is random by default but
can be chosen independently, without a key-derived binding to PeerID. Transport authentication
therefore does not validate DHT placement or provide Sybil resistance.
Sources: dht/routing.py:34--64,111; dht/node.py:105,171,411--462;
dht/protocol.py:149--158; p2p\_daemon\_bindings/control.py:370.}
\end{frame}


% 7 --------------------------------------------------------------------------
\begin{chalmerscover}
\begin{visual}
  \node[text=white,font=\fontsize{32}{36}\selectfont\bfseries] at (7.2,2.8) {Questions?};
\end{visual}
\note{Pause for questions on the architecture before continuing with project progress
and the experimental findings.}
\end{chalmerscover}

\section{Attack surfaces}
% 8 --------------------------------------------------------------------------
\begin{frame}[label=result-custody]{1. Become the record holders}
\begin{visual}
  \node[box,draw=none,minimum width=5.8cm,minimum height=1.1cm,
    inner sep=0pt] at (3.5,4.95) {};
  \node[font=\small\bfseries,text=ChalmersPurple] at (3.5,5.13) {Known group discovery key};
  \node[font=\footnotesize\ttfamily,text=ChalmersPurple] at (3.5,4.77)
    {<prefix>.0b<group\_bits>};
  \node[font=\Large,text=muted] at (7.2,4.95) {$+$};
  \node[box,draw=none,minimum width=5.8cm,minimum height=1.1cm,inner sep=0pt,
    font=\small\bfseries,text=ChalmersPurple] at (10.9,4.95) {Self-chosen \texttt{node\_id}};

  % Leave a clear gap below the premise row and between the two lower blocks.
  \node[box,fill=white,minimum width=3.9cm,minimum height=1.3cm]
    (advertisements) at (2.65,2.4) {Group advertisements\\[2mm]$K$ = \texttt{my\_job.0b}};
  \node[quiet] at (2.65,1.1) {Empty group bits by default};
  \draw[control] (advertisements.east)--node[above,font=\footnotesize]{stored on}(7.55,2.4);
  \draw[ctl!50,dashed,rounded corners=4mm,fill=ctl!3] (7.75,1.55) rectangle (14.0,3.25);
  \node[heading,text=ctl] at (10.875,3.65) {5 nodes nearest to $K$};
  \foreach \xx in {8.5,9.7,10.9,12.1,13.3}
    \node[attacker,minimum size=9mm] at (\xx,2.4) {A};
  \node[quiet] at (10.875,1.1) {Sybils with crafted \texttt{node\_id} values};
  \node[labeltext,text=ChalmersPurple] at (7.2,0.25)
    {All five replicas sit on attacker-run nodes.};
\end{visual}
\resultlink{table-custody_live}{Exact tables}
\note{Ask: what if all five record holders are ours? The group discovery key has format
\texttt{<prefix>.0b<group\_bits>}. Prefix names the averaging namespace; group\_bits is an
optional binary shard suffix. Default empty bits give my\_job.0b, fixed for that prefix.
A sharded example is my\_job.0b010; current co-members can predict the next shard.
This discovery key is different from the random assembled group ID.
\par Attackers independently choose node\_id values at XOR distances 0--4 from the key's
hash and introduce these nodes into the overlay. With five replicas by default, all five
holders are attacker-run. This is measured custody of advertisements under this key;
honest cached copies can still exist. DHT-only Sybils need not train.
\par Placement and reply filtering are stages of one attack. The next slide shows tailored
replies. Five holders alone do not guarantee isolation; the later lookup experiment explains
the additional coverage needed.
Sources: HANDOFF 4.1,4.5--4.7; key\_manager.py:42--44;
step5\_selective.py:158--161; dht/node.py:105,411--462.}
\end{frame}

% 9 --------------------------------------------------------------------------
\begin{frame}[label=result-selective]{2. Give the victim a different view}
\begin{visual}
  % The replica placement was established on slide 8; focus here on the replies.
  \node[box,draw=ctl!35,fill=ctl!4,minimum width=6.2cm,minimum height=0.8cm,
    font=\small\bfseries,text=ctl] (replicas) at (7.2,4.85) {Attacker-run DHT replicas};
  \draw[control] (5.2,4.4)--(3.5,3.28);
  \draw[flow] (9.2,4.4)--(10.9,3.28);

  \fill[ctl!4,rounded corners=3mm] (0.55,0.45) rectangle (6.45,3.15);
  \fill[ChalmersPurple!4,rounded corners=3mm] (7.95,0.45) rectangle (13.85,3.15);
  \node[quiet] at (3.5,2.85) {Reply to};
  \node[quiet] at (10.9,2.85) {Reply to};
  \node[heading] at (3.5,2.43) {the victim};
  \node[heading] at (10.9,2.43) {other honest readers};

  \node[attacker,minimum size=8mm] at (3.5,1.61) {A};
  \node[peer,minimum size=7mm] at (9.65,1.61) {H};
  \node[peer,minimum size=7mm] at (10.9,1.61) {H};
  \node[attacker,minimum size=8mm] at (12.15,1.61) {A};
  \node[quiet,text=ctl] at (3.5,0.85) {Attacker candidates only};
  \node[quiet,text=ChalmersPurple] at (10.9,0.85) {Unfiltered candidates};
  \node[labeltext,text=ChalmersPurple] at (7.2,0.0)
    {Goal: steer the victim into an attacker-led group.};
\end{visual}
\resultlink{table-selective}{Exact tables}
\note{Goal: steer the victim toward an attacker-led group, matching step5\_selective.
Holders accept honest advertisements. Each responder checks the authenticated requester's
PeerID: return the full dictionary to other readers, but retain only attacker advertisements
for the victim. Symbols identify advertised candidates, not a formed group. Unfiltered
replies include both honest (H) and attacker (A) advertisements.
\par In the cold-key selective run, all 12 honest writes succeeded. The victim saw zero
honest candidates and one attacker, while the bystander saw all 12 honest candidates and
the attacker. This demonstrates selective discovery, not guaranteed admission or isolation.
Honest replies can merge omitted entries back; in the warm-key control they did so.
\par The later sustained-capture result uses a different policy: empty replies to the victim,
hide the victim from other honest readers, and truthful replies to attackers, who join the
victim. Twenty crafted identities covered the default lookup candidates; five replica holders
alone were insufficient. Sources: HANDOFF 4.5--4.7; step5\_selective.py:84--111;
results/step5\_selective.csv.}
\end{frame}

\section{Findings}
% 10 --------------------------------------------------------------------------
\begin{frame}[label=result-admission]{$\beta$ is chosen, not inherited}
\begin{visual}
  \node[labeltext,text=ChalmersPurple] at (7.2,5.1)
    {When the attacker leads, it admits only the victim + its own peers.};
  \foreach \xx/\title/\k/\share/\size in {
    2.45/{1 attacker}/1/0.500/2,
    7.2/{2 attackers}/2/0.667/3,
    11.95/{4 attackers}/4/0.800/5} {
    \fill[ChalmersPurple!4,rounded corners=3mm]
      (\xx-2.075,1.25) rectangle (\xx+2.075,4.5);
    \node[heading] at (\xx,4.06) {\title};
    \node[quiet] at (\xx,3.58) {$\alpha=\k/16$ of the swarm};
    \node[font=\Large,text=ctl] at (\xx,2.05) {$\beta=\share$};
    \node[quiet] at (\xx,1.57) {\size\ peers in the group};
  }

  \node[victim,minimum size=6mm] at (2.0,2.87) {V};
  \node[attacker,minimum size=6mm] at (2.9,2.87) {A};
  \node[victim,minimum size=6mm] at (6.25,2.87) {V};
  \foreach \xx in {7.2,8.15}
    \node[attacker,minimum size=6mm] at (\xx,2.87) {A};
  \node[victim,minimum size=6mm] at (10.35,2.87) {V};
  \foreach \xx in {11.15,11.95,12.75,13.55}
    \node[attacker,minimum size=6mm] at (\xx,2.87) {A};

  \node[quiet] at (7.2,0.79)
    {Target group size: 8\quad\textbar\quad minimum: 2 (default)};
  \node[labeltext,text=ChalmersPurple] at (7.2,0.17)
    {Even 50\% attacker share is outside Krum's guarantee.};
\end{visual}
\resultlink{table-admission}{Exact tables}
\note{Transition: ``The goal is to steer the victim into an attacker-led group.
When the attacker does lead, how much control does it have?''
\par Setup: live local swarms, toy regression, unchanged core. Admission was measured
separately from discovery filtering. All peers targeted eight members; the minimum remained
two. The leader rejects other honest peers and can close below target at its deadline.
\par Alpha is the swarm's attacker fraction; beta is the group's. These observed groups
contain the victim plus k attackers, so beta=k/(k+1). Partially filled groups also occurred:
the examples do not imply every attacker-led round includes all k attackers.
For k=1, pairs occurred in 12/21 victim rounds. Beta=0.5 is 8 times alpha=1/16;
mean beta=0.315 is about 5 times alpha, not 8.
\par Krum requires 2f+2<n, with f Byzantine members among n. Even n=2,f=1 violates
that condition. This removes its guarantee; Krum was not evaluated here.
Sources: step3\_capture.py:138--159,220--250; averager.py:119; matchmaking.py:295--297;
step3 capture CSVs d1.0 k1,2,4;
\href{https://papers.nips.cc/paper/2017/file/f4b9ec30ad9f68f89b29639786cb62ef-Paper.pdf}{Blanchard et al., NIPS 2017}.}
\end{frame}

% 11 -------------------------------------------------------------------------
\begin{frame}[label=result-lookup]{Caching}
\begin{visual}
  \node[quiet] at (7.2,5.33)
    {Lookup starts with up to 20 nearest known peers (default)};

%   % One arrow summarises the holders' selective replies for the same key.
  \fill[ctl!5,rounded corners=3mm] (0.45,3.2) rectangle (4.05,4.88);
  \node[heading,text=ctl] at (2.25,4.53) {5 Sybil holders};
  \foreach \xx in {0.95,1.6,2.25,2.9,3.55}
    \node[attacker,minimum size=5.5mm,font=\scriptsize\bfseries] at (\xx,3.83) {S};
  \draw[control] (4.05,3.83)--(9.5,3.83);
  \node[quiet,text=ctl] at (6.75,4.42) {Filtered replies};
  \node[box,fill=white,draw=ctl!35,minimum width=3.2cm,minimum height=7mm,
    inner sep=1mm,font=\small\bfseries,text=ctl] at (6.75,3.83) {A};

  \fill[ChalmersPurple!5,rounded corners=3mm] (0.45,1.35) rectangle (4.05,2.95);
  \node[heading] at (2.25,2.6) {Honest cache};
  \node[peer,minimum size=7mm] at (2.25,1.95) {H};
  \draw[flow] (4.05,1.95)--(9.5,1.95);
  \node[quiet] at (6.75,2.54) {Cached reply};
  \node[box,fill=white,minimum width=3.2cm,minimum height=7mm,
    inner sep=1mm,text=ChalmersPurple] at (6.75,1.95)
    {H1, H2, H3, \textcolor{ctl}{A}};

  \fill[ChalmersPurple!4,rounded corners=3mm] (9.5,1.35) rectangle (14.1,4.88);
  \node[heading] at (11.8,4.53) {Victim's lookup};
  \node[victim] (reader) at (11.8,3.8) {V};
  \draw[flow] (reader.south)--(11.8,3.0);
  \node[box,fill=white,minimum width=4cm,minimum height=1.05cm,
    inner sep=1mm] at (11.8,2.47)
    {\footnotesize Merged candidates\\[2mm]
     \textcolor{ChalmersPurple}{H1, H2, H3,} \textcolor{ctl}{\textbf{A}}};
  \node[quiet,text=ChalmersPurple] at (11.8,1.64) {Honest candidates remain};

  \draw[muted!30] (0.45,1.08)--(14.1,1.08);
  \node[heading,text=ChalmersPurple] at (3.5,0.7) {5 Sybils: 0\% capture};
  \node[heading,text=ctl] at (10.9,0.7) {20 Sybils: 100\% capture};
  \node[quiet] at (7.2,0.14)
    {Capture rates: separate empty-reply experiment\quad\textbar\quad caching on};
\end{visual}
\resultlink{table-lookup_threshold}{Exact tables}
\note{Replicas store PeerID advertisements: availability and expiry, not gradients.
S denotes DHT-only Sybil holders; A is an advertised training attacker. Arrows show replies
for one key. Selective filtering returns only A, while an honest cache restores H1--H3.
The result is a candidate list, not a formed group.
\par Discovery uses latest=True and merges dictionary entries by PeerID: repeated replies
advertising A still create one entry. A remains eligible, but selection follows deadline
and PeerID order, not a one-in-four draw. An early deadline can favour A; honest entries
prevent complete control of discovery. Existing local entries also survive until expiry.
\par Rates below are separate chase runs: empty victim replies, attackers join the victim.
8 training peers (3 attackers), additional DHT Sybils, caching on. Five Sybils: 0/15 and
0/15 captured rounds; twenty: 11/11 and 9/9. Capture means victim plus attackers only.
At 18, mean capture was 14 percent; bucket\_size=40: 25 Sybils gave 9 percent, 45 gave
100 percent. Sources: HANDOFF 4.5--4.6; step5\_selective.py:100--107;
matchmaking.py:478--481; dht/node.py:834--838; step6 chase cached CSVs.}
\end{frame}

% 12 -------------------------------------------------------------------------
\begin{frame}[label=result-sustained]{Capture persisted across rounds}
\begin{visual}
  \node[labeltext] at (7.2,5.15)
    {8 training peers (3 attackers)\quad+\quad25 DHT-only Sybils};
  \node[heading,text=ChalmersPurple] at (3.0,4.5) {Run};
  \node[heading,text=ChalmersPurple] at (7.2,4.5) {Captured rounds};
  \node[heading,text=ChalmersPurple] at (11.4,4.5) {Recorded rounds};
  \draw[ChalmersPurple!50,line width=0.7pt] (1.0,4.17)--(13.4,4.17);
  \foreach \yy/\runindex/\roundtotal in {3.72/1/15,2.97/2/13,2.22/3/12} {
    \node[font=\large] at (3.0,\yy) {\runindex};
    \node[font=\large\bfseries] at (7.2,\yy) {\roundtotal};
    \node[font=\large] at (11.4,\yy) {\roundtotal};
    \draw[ChalmersPurple!15] (1.0,\yy-0.35)--(13.4,\yy-0.35);
  }
  \node[heading,text=ctl] at (7.2,1.44)
    {All 40 recorded rounds: victim + attacker peers only.};
  \node[quiet] at (7.2,0.86)
    {Sybils joined before training\quad\textbar\quad caching on\quad\textbar\quad sharding off};
  \draw[ChalmersPurple!25] (0.6,0.5)--(13.8,0.5);
  \node[font=\small\bfseries,text=ChalmersPurple] at (2.25,0.13) {16 training peers};
  \node[quiet,text=ink] at (7.0,0.13) {3 attempts\quad\textbar\quad0 victim rounds};
  \node[quiet,text=ChalmersPurple] at (11.9,0.13) {Capture not assessable};
  \node[quiet] at (7.2,-0.43)
    {Arrival-to-capture delay unmeasured\quad\textbar\quad16-peer attempts: handoff report};
\end{visual}
\resultlink{table-sustained}{Exact tables}
\note{Did capture persist when the victim averaged again? Each table row is a separate run.
Captured means the victim averaged with attacker peers and no other honest peer.
The captured count equals the recorded count in every run: 15/15,13/13,12/12---40/40.
Each run trained for 15 epochs; epochs and averaging rounds are different counters,
so the runs contain different numbers of recorded averaging rounds.
\par These results use eight training peers, including three attackers, plus 25 DHT-only
Sybils. Caching was on; sharding was off. Sybils joined before training. This establishes
sustained capture after preparation, not arrival-to-capture latency.
\par We also tried sixteen training peers. HANDOFF's configuration notes report three
attempts with target group size explicitly set: only 4--6 peers ever averaged, and the
victim had zero recorded rounds. There is no denominator for capture, so this is not a
zero-percent result or evidence of protection at larger scale. Matching raw outputs were
not located; the separate strip is a handoff report, not another measured capture rate.
\par Sources: HANDOFF 4.6--4.7 and configuration traps (lines 690--692);
step6 cached s25 repeats 1--3. Mechanism and timing details: PRESENTATION.md.}
\end{frame}

% 13 -------------------------------------------------------------------------
\begin{frame}[label=result-tracking]{Tracking the bucket, chasing the group}
\begin{visual}
  \node[labeltext] at (7.2,5.15)
    {Victim rounds with at least one attacker in the group};
  \node[heading,text=ChalmersPurple] at (3.0,4.5) {Training peers};
  \node[heading,text=ChalmersPurple] at (7.2,4.5) {Fixed bucket};
  \node[heading,text=ChalmersPurple] at (11.4,4.5) {Tracks victim};
  \draw[ChalmersPurple!50,line width=0.7pt] (1.0,4.17)--(13.4,4.17);
  \foreach \yy/\population/\fixedrate/\trackingrate in {
    3.72/16/56/84,2.97/32/40/72,2.22/48/18/38} {
    \node[font=\large] at (3.0,\yy) {\population};
    \node[font=\large] at (7.2,\yy) {\fixedrate\%};
    \node[font=\large\bfseries] at (11.4,\yy) {\trackingrate\%};
    \draw[ChalmersPurple!15] (1.0,\yy-0.35)--(13.4,\yy-0.35);
  }
  \node[labeltext] at (7.2,1.34)
    {A correct bucket prediction does not guarantee the same averaging group.};
  \node[quiet] at (7.2,0.72)
    {Next-bucket predictions: 63/63 correct, using shared group information.};
  \node[quiet] at (7.2,0.02)
    {Sharding on\quad\textbar\quad4 discovery buckets\quad\textbar\quad4 attackers\quad\textbar\quad2 runs per setting};
\end{visual}
\resultlink{table-tracking}{Exact tables}
\note{Here we enable optional sharding: peers use four discovery buckets to find averaging
partners. The default has sharding off. An attacker first needs to share the victim's
group. It receives the group\_id and ordered member list, calculates the victim's next
bucket using the victim's index, and overrides its own assignment to follow. If they
share another group, it gets fresh information; otherwise it cannot predict the victim's
next move until they meet again. Attackers operate independently.
\par The table measures victim rounds containing at least one attacker, not full capture.
Other honest peers may remain. Fixed-bucket attackers stay in one bucket; tracking
attackers follow when they can. Both arms have four attacker identities. Tracking helps
at every size, but as buckets contain more groups, choosing the right bucket is less
likely to produce the same group. Rates average two repeats; target group size is four.
\par All 63 prediction records were correct across six runs and 49 distinct shared groups.
Several attackers can predict the same transition. This does not establish permanent
tracking. Tracking and selective admission were not measured together.
Sources: HANDOFF 4.4; step4 CSVs; key\_manager.py:94--104; step4\_reformation.py:183--191.}
\end{frame}

% 14 -------------------------------------------------------------------------
\begin{frame}[label=result-poisoning]{Poisoning: results and limits}
\begin{visual}
  \node[labeltext] at (7.2,5.15)
    {8 training peers\quad\textbar\quad1 poisoner\quad\textbar\quad64-parameter regression};
  \foreach \yy/\statement in {
    4.0/{One peer reversed and scaled its outgoing gradients, then claimed an averaging weight of one million.},
    2.85/{Weighted mean left honest peers farther from the optimum than in matching clean runs.},
    1.7/{Coordinate median substantially reduced the damage across all three seeds.}} {
    \fill[ChalmersPurple] (0.75,\yy) circle (1.1mm);
    \node[anchor=west,text width=12.0cm,align=left,font=\normalsize]
      at (1.1,\yy) {\statement};
  }
  \draw[ChalmersPurple!25] (0.6,0.83)--(13.8,0.83);
  \node[quiet] at (7.2,0.43)
    {Poisoning was tested separately from group capture.};
  \node[quiet] at (7.2,0.02)
    {Krum: theoretical limits only; no implementation or experiment.};
\end{visual}
\resultlink{table-poisoning}{Exact tables}
\note{This separate learning experiment used eight peers and one poisoner on a small
regression model. The attacker multiplied outgoing gradient slices by minus five and
claimed an averaging weight of one million. Normally that weight reflects the sample
count. The inflated claim made its bad gradient dominate the weighted average. Honest
peers ended farther from the optimum than in matching clean runs. Coordinate median
ignores declared weights and substantially reduced the damage.
\par Poisoning and group capture were tested separately. Krum was considered
theoretically, not implemented or tested. This is toy-model evidence, not transformer evidence.
\par Metric: honest peers' mean Euclidean distance at epoch 20, divided by the same-seed
clean weighted-mean control. Weighted poison: 45.69, 148.48, 35.01 times control;
median: 0.92, 2.32, 0.50. Ordinary-weight sign flip has one seed: 0.80 times control.
Target group size was four; actual membership was not recorded, so no fixed local
beta or damage-versus-beta curve was established.
Sources: HANDOFF 4.11; step7\_poison CSVs; exact appendix tables.}
\end{frame}

\section{Limitations}
% 15 -------------------------------------------------------------------------
\begin{frame}[label=result-limitations]{Limitations}
\begin{visual}
  \node[anchor=west,font=\small\bfseries,text=ChalmersPurple] at (2.65,5.23)
    {Tested setting};
  \node[anchor=west,font=\small\bfseries,text=ChalmersPurple] at (8.1,5.23)
    {Limit of the evidence};
  \draw[ChalmersPurple!50,line width=0.7pt] (0.45,4.86)--(14.0,4.86);
  \foreach \yy/\topic/\testedsetting/\evidencelimit in {
    4.38/{Model}/{Linear regression\\64 parameters}/{No transformer experiments},
    3.51/{Network}/{8--104 nodes on one machine}/{Large-network routing untested},
    2.64/{Timing}/{Sybils joined before training}/{Time to establish capture\\not measured},
    1.77/{Attack}/{Stages tested separately}/{Full attack chain unmeasured},
    0.90/{Krum}/{Theoretical limits only\\Not implemented}/{Needs full-gradient distances;\\reducers receive slices}} {
    \node[anchor=west,font=\small\bfseries,text=ChalmersPurple,
      text width=2cm,align=left] at (0.45,\yy) {\topic};
    \node[anchor=west,font=\small,text width=4.8cm,align=left]
      at (2.65,\yy) {\testedsetting};
    \node[anchor=west,font=\small,text width=5.6cm,align=left]
      at (8.1,\yy) {\evidencelimit};
    \draw[ChalmersPurple!15] (0.45,\yy-0.43)--(14.0,\yy-0.43);
  }
  \node[quiet,anchor=west,align=left] at (0.45,0.02)
    {Control-plane defences are untested; default group mixing remains unsettled.};
\end{visual}
\resultlink{table-index}{All experiment tables}
\note{The toy regression is not transformer evidence. All swarms ran on one machine;
idle fillers at 104 nodes did not expand the victim's routing view. Capture followed
Sybil preparation. Default mixing differed at 16 and 24 peers; load was not controlled.
The full attack chain remains unmeasured; poisoning did not record actual group composition.
\par Krum is absent from the audited Hivemind and was not implemented in this project.
I only checked its theoretical condition, 2f+2<n. Violating it removes the guarantee;
it is not an experimental demonstration of Krum failing.
\par Krum compares complete gradients using pairwise distances. A straightforward
implementation collects the full gradient matrix (one vector per peer) at a trusted
central coordinator. Hivemind's reducers receive assigned slices; trainers retain their
own full gradients. Krum cannot directly replace this slice reducer. A central server
is one option, not a requirement: distributed distances and consistent vector selection
would need extra communication and coordination. Neither approach was implemented.
\par Coordinate median WAS evaluated; control-plane defences remain untested.
Sources: HANDOFF 4.6--4.11; step7\_poison.py; partition.py; Blanchard et al., NeurIPS 2017.}
\end{frame}

% \appendix
% \renewcommand{\Chalmersinsertpagenumber}{A\insertframenumberinappendix\,/\,\insertappendixframenumber}
% \section{Experiment tables}
% % Generated by ../260918_Project_Presentation_support/scripts/render_experiment_tables.py.
% Edit that renderer or the support JSON to regenerate these tables.
% Native Beamer appendix; parent supplies array, visual, colours and resultlink.
\begin{frame}[label=table-index]{Recorded experiment tables}
\begin{visual}
  \node[quiet] at (7.2,5.28) {Saved output, exact counts and six-decimal summaries};
  \node[anchor=west,draw=ChalmersPurple!20,fill=ChalmersPurple!3,rounded corners=1mm,inner sep=1.8mm,font=\footnotesize,text=ChalmersPurple] at (0.45,4.72) {\hyperlink{table-custody}{\makebox[6.0cm][l]{01\quad Record custody: placement}}};
  \node[anchor=west,draw=ChalmersPurple!20,fill=ChalmersPurple!3,rounded corners=1mm,inner sep=1.8mm,font=\footnotesize,text=ChalmersPurple] at (0.45,4.23) {\hyperlink{table-custody_live}{\makebox[6.0cm][l]{02\quad Record custody: live check}}};
  \node[anchor=west,draw=ChalmersPurple!20,fill=ChalmersPurple!3,rounded corners=1mm,inner sep=1.8mm,font=\footnotesize,text=ChalmersPurple] at (0.45,3.74) {\hyperlink{table-selective}{\makebox[6.0cm][l]{03\quad Selective discovery}}};
  \node[anchor=west,draw=ChalmersPurple!20,fill=ChalmersPurple!3,rounded corners=1mm,inner sep=1.8mm,font=\footnotesize,text=ChalmersPurple] at (0.45,3.25) {\hyperlink{table-warm_records}{\makebox[6.0cm][l]{04\quad Warm records and expiry}}};
  \node[anchor=west,draw=ChalmersPurple!20,fill=ChalmersPurple!3,rounded corners=1mm,inner sep=1.8mm,font=\footnotesize,text=ChalmersPurple] at (0.45,2.76) {\hyperlink{table-candidate_flood}{\makebox[6.0cm][l]{05\quad Candidate flooding}}};
  \node[anchor=west,draw=ChalmersPurple!20,fill=ChalmersPurple!3,rounded corners=1mm,inner sep=1.8mm,font=\footnotesize,text=ChalmersPurple] at (0.45,2.27) {\hyperlink{table-admission}{\makebox[6.0cm][l]{06\quad Admission and beta}}};
  \node[anchor=west,draw=ChalmersPurple!20,fill=ChalmersPurple!3,rounded corners=1mm,inner sep=1.8mm,font=\footnotesize,text=ChalmersPurple] at (0.45,1.78) {\hyperlink{table-lookup_threshold}{\makebox[6.0cm][l]{07\quad Lookup coverage threshold}}};
  \node[anchor=west,draw=ChalmersPurple!20,fill=ChalmersPurple!3,rounded corners=1mm,inner sep=1.8mm,font=\footnotesize,text=ChalmersPurple] at (0.45,1.29) {\hyperlink{table-lookup_controls}{\makebox[6.0cm][l]{08\quad Lookup controls}}};
  \node[anchor=west,draw=ChalmersPurple!20,fill=ChalmersPurple!3,rounded corners=1mm,inner sep=1.8mm,font=\footnotesize,text=ChalmersPurple] at (0.45,0.80) {\hyperlink{table-lookup_history}{\makebox[6.0cm][l]{09\quad Earlier lookup controls}}};
  \node[anchor=west,draw=ChalmersPurple!20,fill=ChalmersPurple!3,rounded corners=1mm,inner sep=1.8mm,font=\footnotesize,text=ChalmersPurple] at (7.6,4.72) {\hyperlink{table-sustained}{\makebox[6.0cm][l]{10\quad Sustained capture}}};
  \node[anchor=west,draw=ChalmersPurple!20,fill=ChalmersPurple!3,rounded corners=1mm,inner sep=1.8mm,font=\footnotesize,text=ChalmersPurple] at (7.6,4.23) {\hyperlink{table-tracking}{\makebox[6.0cm][l]{11\quad Re-formation and tracking}}};
  \node[anchor=west,draw=ChalmersPurple!20,fill=ChalmersPurple!3,rounded corners=1mm,inner sep=1.8mm,font=\footnotesize,text=ChalmersPurple] at (7.6,3.74) {\hyperlink{table-poisoning}{\makebox[6.0cm][l]{12\quad Poisoning: endpoint distances}}};
  \node[anchor=west,draw=ChalmersPurple!20,fill=ChalmersPurple!3,rounded corners=1mm,inner sep=1.8mm,font=\footnotesize,text=ChalmersPurple] at (7.6,3.25) {\hyperlink{table-poisoning_ratios}{\makebox[6.0cm][l]{13\quad Poisoning: control ratios}}};
  \node[anchor=west,draw=ChalmersPurple!20,fill=ChalmersPurple!3,rounded corners=1mm,inner sep=1.8mm,font=\footnotesize,text=ChalmersPurple] at (7.6,2.76) {\hyperlink{table-withdrawal}{\makebox[6.0cm][l]{14\quad Withdrawal: negative capture result}}};
  \node[anchor=west,draw=ChalmersPurple!20,fill=ChalmersPurple!3,rounded corners=1mm,inner sep=1.8mm,font=\footnotesize,text=ChalmersPurple] at (7.6,2.27) {\hyperlink{table-withdrawal_discovery}{\makebox[6.0cm][l]{15\quad Withdrawal: DHT visibility}}};
  \node[anchor=west,draw=ChalmersPurple!20,fill=ChalmersPurple!3,rounded corners=1mm,inner sep=1.8mm,font=\footnotesize,text=ChalmersPurple] at (7.6,1.78) {\hyperlink{table-baseline}{\makebox[6.0cm][l]{16\quad Baseline training}}};
  \node[anchor=west,draw=ChalmersPurple!20,fill=ChalmersPurple!3,rounded corners=1mm,inner sep=1.8mm,font=\footnotesize,text=ChalmersPurple] at (7.6,1.29) {\hyperlink{table-mixing}{\makebox[6.0cm][l]{17\quad Baseline membership mixing}}};
  \node[anchor=west,draw=ChalmersPurple!20,fill=ChalmersPurple!3,rounded corners=1mm,inner sep=1.8mm,font=\footnotesize,text=ChalmersPurple] at (7.6,0.80) {\hyperlink{table-sharding_handoff}{\makebox[6.0cm][l]{18\quad Sharding failures (reported only)}}};
  \node[quiet,font=\scriptsize] at (7.2,0.0) {Each table links back to its result slide. Reported-only material is labelled.};
\end{visual}
\resultlink{result-limitations}{Back to main presentation}
\note{Tables are generated from the saved CSV measurements by export\_experiment\_tables.py and rendered by render\_experiment\_tables.py. No experiments were rerun. Decimal summaries are rounded to six places; integer fractions are exact counts. Full paths, definitions and unrounded calculated metrics are in EXPERIMENT\_TABLES.md and experiment-tables-data.json. Sharding-failure observations are explicitly labelled as reported in HANDOFF without matching raw output.}
\end{frame}

\begin{frame}[label=table-custody]{Record custody: placement calculation}
\begin{visual}
  \node[quiet,text=muted,text width=13.5cm,font=\footnotesize] at (7.2,5.3) {5 replica positions; 40 arithmetic repeats per arm and swarm size};
  \node[anchor=north west,inner sep=0pt,font=\footnotesize] at (0.45,4.66) {%
    \begingroup\fontsize{8}{10}\selectfont\setlength{\tabcolsep}{3pt}\renewcommand{\arraystretch}{1.0}%
    \setbox\strutbox=\hbox{\vrule height7pt depth3pt width0pt}%
    \begin{tabular}{>{\raggedright\arraybackslash}p{2.05cm}>{\raggedright\arraybackslash}p{1.65cm}>{\raggedright\arraybackslash}p{2.45cm}>{\raggedright\arraybackslash}p{2.85cm}>{\raggedright\arraybackslash}p{2.85cm}}
      \textcolor{ChalmersPurple}{\textbf{Honest IDs}} & \textcolor{ChalmersPurple}{\textbf{Added IDs}} & \textcolor{ChalmersPurple}{\textbf{Repeats}} & \textcolor{ChalmersPurple}{\textbf{Crafted share}} & \textcolor{ChalmersPurple}{\textbf{Random share}} \\[1mm]\hline
      50 & 5 & 40 per arm & 100.0\% & 9.0\% \\
      1000 & 5 & 40 per arm & 100.0\% & 1.0\% \\
      10000 & 5 & 40 per arm & 100.0\% & 0.0\% \\
    \end{tabular}\endgroup
  };
  \node[quiet,font=\scriptsize,text width=13.3cm] at (7.2,0.65) {Share of replica positions; not a live group-capture experiment.};
  \node[font=\scriptsize] at (5.6,0.15) {\hyperlink{table-index}{\beamerbutton{All tables}}};
  \node[font=\scriptsize] at (8.6,0.15) {\hyperlink{table-custody_live}{\beamerbutton{Next table}}};
\end{visual}
\resultlink{result-custody}{Back to result}
\note{Share of the five replica positions, not group capture. Arithmetic experiment; no live swarm at these sizes. Source: research/hivemind/results/step3\_eclipse\_analytical.csv. Every displayed crafted repeat has all five replicas. Exact saved counts; displayed decimals are rounded. Source paths and full per-run calculations: EXPERIMENT\_TABLES.md and experiment-tables-data.json; regenerated by export\_experiment\_tables.py.}
\end{frame}

\begin{frame}[label=table-custody_live]{Record custody: live DHT nodes}
\begin{visual}
  \node[quiet,text=muted,text width=13.5cm,font=\footnotesize] at (7.2,5.3) {crafted IDs | 24 honest DHT nodes; three repeats per row; mean share and startup time};
  \node[anchor=north west,inner sep=0pt,font=\scriptsize] at (0.45,4.66) {%
    \begingroup\fontsize{8}{10}\selectfont\setlength{\tabcolsep}{3pt}\renewcommand{\arraystretch}{1.0}%
    \setbox\strutbox=\hbox{\vrule height7pt depth3pt width0pt}%
    \begin{tabular}{>{\raggedright\arraybackslash}p{1.6cm}>{\raggedright\arraybackslash}p{1.6cm}>{\raggedright\arraybackslash}p{1.15cm}>{\raggedright\arraybackslash}p{2.6cm}>{\raggedright\arraybackslash}p{2.5cm}>{\raggedright\arraybackslash}p{2.35cm}}
      \textcolor{ChalmersPurple}{\textbf{ID choice}} & \textcolor{ChalmersPurple}{\textbf{Added nodes}} & \textcolor{ChalmersPurple}{\textbf{Repeats}} & \textcolor{ChalmersPurple}{\textbf{Predicted share}} & \textcolor{ChalmersPurple}{\textbf{Stored share}} & \textcolor{ChalmersPurple}{\textbf{Build time (s)}} \\[1mm]\hline
      crafted & 0 & 3 & 0.000000 & 0.000000 & 0.000000 \\
      crafted & 1 & 3 & 0.200000 & 0.200000 & 0.100667 \\
      crafted & 3 & 3 & 0.600000 & 0.600000 & 0.299667 \\
      crafted & 5 & 3 & 1.000000 & 1.000000 & 0.505333 \\
      crafted & 8 & 3 & 1.000000 & 1.000000 & 0.796667 \\
    \end{tabular}\endgroup
  };
  \node[quiet,font=\scriptsize,text width=13.3cm] at (7.2,0.65) {All 30 individual observations match predicted custody. Startup is not interception latency.};
  \node[font=\scriptsize] at (5.6,0.15) {\hyperlink{table-index}{\beamerbutton{All tables}}};
  \node[font=\scriptsize] at (8.6,0.15) {\hyperlink{table-custody_live-2}{\beamerbutton{Next table}}};
\end{visual}
\resultlink{result-custody}{Back to result}
\note{24 honest DHT nodes. Replica share, not group capture. Build time includes process startup; it is not interception latency. All 30 CSV rows have structural\_capture equal to store\_capture. Source: research/hivemind/results/step3\_eclipse\_live.csv. Exact saved counts; displayed decimals are rounded. Source paths and full per-run calculations: EXPERIMENT\_TABLES.md and experiment-tables-data.json; regenerated by export\_experiment\_tables.py.}
\end{frame}

\begin{frame}[label=table-custody_live-2]{Record custody: live DHT nodes}
\begin{visual}
  \node[quiet,text=muted,text width=13.5cm,font=\footnotesize] at (7.2,5.3) {random IDs | 24 honest DHT nodes; three repeats per row; mean share and startup time};
  \node[anchor=north west,inner sep=0pt,font=\scriptsize] at (0.45,4.66) {%
    \begingroup\fontsize{8}{10}\selectfont\setlength{\tabcolsep}{3pt}\renewcommand{\arraystretch}{1.0}%
    \setbox\strutbox=\hbox{\vrule height7pt depth3pt width0pt}%
    \begin{tabular}{>{\raggedright\arraybackslash}p{1.6cm}>{\raggedright\arraybackslash}p{1.6cm}>{\raggedright\arraybackslash}p{1.15cm}>{\raggedright\arraybackslash}p{2.6cm}>{\raggedright\arraybackslash}p{2.5cm}>{\raggedright\arraybackslash}p{2.35cm}}
      \textcolor{ChalmersPurple}{\textbf{ID choice}} & \textcolor{ChalmersPurple}{\textbf{Added nodes}} & \textcolor{ChalmersPurple}{\textbf{Repeats}} & \textcolor{ChalmersPurple}{\textbf{Predicted share}} & \textcolor{ChalmersPurple}{\textbf{Stored share}} & \textcolor{ChalmersPurple}{\textbf{Build time (s)}} \\[1mm]\hline
      random & 0 & 3 & 0.000000 & 0.000000 & 0.000000 \\
      random & 1 & 3 & 0.000000 & 0.000000 & 0.098000 \\
      random & 3 & 3 & 0.133333 & 0.133333 & 0.296000 \\
      random & 5 & 3 & 0.266667 & 0.266667 & 0.836667 \\
      random & 8 & 3 & 0.466667 & 0.466667 & 0.808333 \\
    \end{tabular}\endgroup
  };
  \node[quiet,font=\scriptsize,text width=13.3cm] at (7.2,0.65) {All 30 individual observations match predicted custody. Startup is not interception latency.};
  \node[font=\scriptsize] at (5.6,0.15) {\hyperlink{table-index}{\beamerbutton{All tables}}};
  \node[font=\scriptsize] at (8.6,0.15) {\hyperlink{table-selective}{\beamerbutton{Next table}}};
\end{visual}
\resultlink{result-custody}{Back to result}
\note{24 honest DHT nodes. Replica share, not group capture. Build time includes process startup; it is not interception latency. All 30 CSV rows have structural\_capture equal to store\_capture. Source: research/hivemind/results/step3\_eclipse\_live.csv. Exact saved counts; displayed decimals are rounded. Source paths and full per-run calculations: EXPERIMENT\_TABLES.md and experiment-tables-data.json; regenerated by export\_experiment\_tables.py.}
\end{frame}

\begin{frame}[label=table-selective]{Selective discovery: the recorded replies}
\begin{visual}
  \node[quiet,text=muted,text width=13.5cm,font=\footnotesize] at (7.2,5.3) {Cold key; 12 honest advertisements; caching off; one run per arm};
  \node[anchor=north west,inner sep=0pt,font=\footnotesize] at (0.45,4.66) {%
    \begingroup\fontsize{8}{10}\selectfont\setlength{\tabcolsep}{3pt}\renewcommand{\arraystretch}{1.0}%
    \setbox\strutbox=\hbox{\vrule height7pt depth3pt width0pt}%
    \begin{tabular}{>{\raggedright\arraybackslash}p{2.3cm}>{\raggedright\arraybackslash}p{3.2cm}>{\raggedright\arraybackslash}p{3.2cm}>{\raggedright\arraybackslash}p{3.2cm}}
      \textcolor{ChalmersPurple}{\textbf{Arm}} & \textcolor{ChalmersPurple}{\textbf{Honest writes accepted}} & \textcolor{ChalmersPurple}{\textbf{Victim sees}} & \textcolor{ChalmersPurple}{\textbf{Bystander sees}} \\[1mm]\hline
      honest & 12/12 & 12 H + 0 A & 12 H + 0 A \\
      censor & 0/12 & 0 H + 1 A & 0 H + 1 A \\
      selective & 12/12 & 0 H + 1 A & 12 H + 1 A \\
    \end{tabular}\endgroup
  };
  \node[quiet,font=\scriptsize,text width=13.3cm] at (7.2,0.65) {H = honest candidate; A = attacker candidate. These are advertisements, not formed groups.};
  \node[font=\scriptsize] at (5.6,0.15) {\hyperlink{table-index}{\beamerbutton{All tables}}};
  \node[font=\scriptsize] at (8.6,0.15) {\hyperlink{table-warm_records}{\beamerbutton{Next table}}};
\end{visual}
\resultlink{result-selective}{Back to result}
\note{12 honest advertisements; attack arms have five crafted DHT nodes. Caching off; one recorded run per arm. Candidate visibility, not formed-group membership. A denotes one advertised attacker. Source: research/hivemind/results/step5\_selective.csv. Exact saved counts; displayed decimals are rounded. Source paths and full per-run calculations: EXPERIMENT\_TABLES.md and experiment-tables-data.json; regenerated by export\_experiment\_tables.py.}
\end{frame}

\begin{frame}[label=table-warm_records]{Warm records: honest entries remain}
\begin{visual}
  \node[quiet,text=muted,text width=13.5cm,font=\footnotesize] at (7.2,5.3) {Honest candidates visible out of 12; one saved run per configuration};
  \node[anchor=north west,inner sep=0pt,font=\footnotesize] at (0.45,4.66) {%
    \begingroup\fontsize{8}{10}\selectfont\setlength{\tabcolsep}{3pt}\renewcommand{\arraystretch}{1.0}%
    \setbox\strutbox=\hbox{\vrule height7pt depth3pt width0pt}%
    \begin{tabular}{>{\raggedright\arraybackslash}p{3.1cm}>{\raggedright\arraybackslash}p{1.6cm}>{\raggedright\arraybackslash}p{3.6cm}>{\raggedright\arraybackslash}p{3.6cm}}
      \textcolor{ChalmersPurple}{\textbf{Key history}} & \textcolor{ChalmersPurple}{\textbf{Cache}} & \textcolor{ChalmersPurple}{\textbf{Censor: honest seen}} & \textcolor{ChalmersPurple}{\textbf{Selective: honest seen}} \\[1mm]\hline
      Cold & off & 0 & 0 \\
      Warm: 1 round & off & -- & 12 \\
      Cold & on & -- & 0 \\
      Warm: 1 round & on & 12 & 12 \\
      Warm: 5 rounds & on & -- & 12 \\
      Warm: expired & on & 0 & 0 \\
    \end{tabular}\endgroup
  };
  \node[quiet,font=\scriptsize,text width=13.3cm] at (7.2,0.65) {Expired: 5 s lifetime, 8 s settle before Sybils join. -- = no saved arm.};
  \node[font=\scriptsize] at (5.6,0.15) {\hyperlink{table-index}{\beamerbutton{All tables}}};
  \node[font=\scriptsize] at (8.6,0.15) {\hyperlink{table-candidate_flood}{\beamerbutton{Next table}}};
\end{visual}
\resultlink{result-lookup}{Back to result}
\note{Counts out of 12 honest advertisements; one saved run for each configuration. -- means no matching saved arm; not imputed from HANDOFF. Expired configuration: lifetime 5 s, settle 8 s, before Sybils are introduced. This does not measure arrival-to-capture delay during training. Exact saved counts; displayed decimals are rounded. Source paths and full per-run calculations: EXPERIMENT\_TABLES.md and experiment-tables-data.json; regenerated by export\_experiment\_tables.py. The visible table omits long CSV filenames; the source table retains one filename per row.}
\end{frame}

\begin{frame}[label=table-candidate_flood]{Flooding versus record custody}
\begin{visual}
  \node[quiet,text=muted,text width=13.5cm,font=\footnotesize] at (7.2,5.3) {8 honest adverts | Cold-key candidate lists; one saved run per row};
  \node[anchor=north west,inner sep=0pt,font=\scriptsize] at (0.45,4.66) {%
    \begingroup\fontsize{8}{10}\selectfont\setlength{\tabcolsep}{3pt}\renewcommand{\arraystretch}{1.0}%
    \setbox\strutbox=\hbox{\vrule height7pt depth3pt width0pt}%
    \begin{tabular}{>{\raggedright\arraybackslash}p{1.3cm}>{\raggedright\arraybackslash}p{1.5cm}>{\raggedright\arraybackslash}p{1.6cm}>{\raggedright\arraybackslash}p{1.5cm}>{\raggedright\arraybackslash}p{1.65cm}>{\raggedright\arraybackslash}p{2.0cm}>{\raggedright\arraybackslash}p{2.15cm}}
      \textcolor{ChalmersPurple}{\textbf{Honest adverts}} & \textcolor{ChalmersPurple}{\textbf{Arm}} & \textcolor{ChalmersPurple}{\textbf{Attacker processes}} & \textcolor{ChalmersPurple}{\textbf{Attacker writes}} & \textcolor{ChalmersPurple}{\textbf{Honest visible}} & \textcolor{ChalmersPurple}{\textbf{Attacker / list}} & \textcolor{ChalmersPurple}{\textbf{Attacker share}} \\[1mm]\hline
      8 & honest & 0 & 0 & 8/8 & 0/8 & 0.000000 \\
      8 & flood & 1 & 8 & 8/8 & 8/16 & 0.500000 \\
      8 & flood & 1 & 32 & 8/8 & 32/40 & 0.800000 \\
      8 & flood & 1 & 128 & 8/8 & 128/136 & 0.941176 \\
      8 & eclipse & 5 & 1 & 0/8 & 1/1 & 1.000000 \\
    \end{tabular}\endgroup
  };
  \node[quiet,font=\scriptsize,text width=13.3cm] at (7.2,0.65) {Flooding adds advertised identities; custody can refuse honest writes. Neither row measures group capture.};
  \node[font=\scriptsize] at (5.6,0.15) {\hyperlink{table-index}{\beamerbutton{All tables}}};
  \node[font=\scriptsize] at (8.6,0.15) {\hyperlink{table-candidate_flood-2}{\beamerbutton{Next table}}};
\end{visual}
\resultlink{result-custody}{Back to result}
\note{All 15 saved rows; one run per configuration. Cold-key discovery, not group capture. Flooding uses one process publishing multiple advertised identities. Eclipse uses five crafted record holders refusing honest advertisements. Source: research/hivemind/results/step3\_candidates\_live.csv. Exact saved counts; displayed decimals are rounded. Source paths and full per-run calculations: EXPERIMENT\_TABLES.md and experiment-tables-data.json; regenerated by export\_experiment\_tables.py.}
\end{frame}

\begin{frame}[label=table-candidate_flood-2]{Flooding versus record custody}
\begin{visual}
  \node[quiet,text=muted,text width=13.5cm,font=\footnotesize] at (7.2,5.3) {16 honest adverts | Cold-key candidate lists; one saved run per row};
  \node[anchor=north west,inner sep=0pt,font=\scriptsize] at (0.45,4.66) {%
    \begingroup\fontsize{8}{10}\selectfont\setlength{\tabcolsep}{3pt}\renewcommand{\arraystretch}{1.0}%
    \setbox\strutbox=\hbox{\vrule height7pt depth3pt width0pt}%
    \begin{tabular}{>{\raggedright\arraybackslash}p{1.3cm}>{\raggedright\arraybackslash}p{1.5cm}>{\raggedright\arraybackslash}p{1.6cm}>{\raggedright\arraybackslash}p{1.5cm}>{\raggedright\arraybackslash}p{1.65cm}>{\raggedright\arraybackslash}p{2.0cm}>{\raggedright\arraybackslash}p{2.15cm}}
      \textcolor{ChalmersPurple}{\textbf{Honest adverts}} & \textcolor{ChalmersPurple}{\textbf{Arm}} & \textcolor{ChalmersPurple}{\textbf{Attacker processes}} & \textcolor{ChalmersPurple}{\textbf{Attacker writes}} & \textcolor{ChalmersPurple}{\textbf{Honest visible}} & \textcolor{ChalmersPurple}{\textbf{Attacker / list}} & \textcolor{ChalmersPurple}{\textbf{Attacker share}} \\[1mm]\hline
      16 & honest & 0 & 0 & 16/16 & 0/16 & 0.000000 \\
      16 & flood & 1 & 8 & 16/16 & 8/24 & 0.333333 \\
      16 & flood & 1 & 32 & 16/16 & 32/48 & 0.666667 \\
      16 & flood & 1 & 128 & 16/16 & 128/144 & 0.888889 \\
      16 & eclipse & 5 & 1 & 0/16 & 1/1 & 1.000000 \\
    \end{tabular}\endgroup
  };
  \node[quiet,font=\scriptsize,text width=13.3cm] at (7.2,0.65) {Flooding adds advertised identities; custody can refuse honest writes. Neither row measures group capture.};
  \node[font=\scriptsize] at (5.6,0.15) {\hyperlink{table-index}{\beamerbutton{All tables}}};
  \node[font=\scriptsize] at (8.6,0.15) {\hyperlink{table-candidate_flood-3}{\beamerbutton{Next table}}};
\end{visual}
\resultlink{result-custody}{Back to result}
\note{All 15 saved rows; one run per configuration. Cold-key discovery, not group capture. Flooding uses one process publishing multiple advertised identities. Eclipse uses five crafted record holders refusing honest advertisements. Source: research/hivemind/results/step3\_candidates\_live.csv. Exact saved counts; displayed decimals are rounded. Source paths and full per-run calculations: EXPERIMENT\_TABLES.md and experiment-tables-data.json; regenerated by export\_experiment\_tables.py.}
\end{frame}

\begin{frame}[label=table-candidate_flood-3]{Flooding versus record custody}
\begin{visual}
  \node[quiet,text=muted,text width=13.5cm,font=\footnotesize] at (7.2,5.3) {32 honest adverts | Cold-key candidate lists; one saved run per row};
  \node[anchor=north west,inner sep=0pt,font=\scriptsize] at (0.45,4.66) {%
    \begingroup\fontsize{8}{10}\selectfont\setlength{\tabcolsep}{3pt}\renewcommand{\arraystretch}{1.0}%
    \setbox\strutbox=\hbox{\vrule height7pt depth3pt width0pt}%
    \begin{tabular}{>{\raggedright\arraybackslash}p{1.3cm}>{\raggedright\arraybackslash}p{1.5cm}>{\raggedright\arraybackslash}p{1.6cm}>{\raggedright\arraybackslash}p{1.5cm}>{\raggedright\arraybackslash}p{1.65cm}>{\raggedright\arraybackslash}p{2.0cm}>{\raggedright\arraybackslash}p{2.15cm}}
      \textcolor{ChalmersPurple}{\textbf{Honest adverts}} & \textcolor{ChalmersPurple}{\textbf{Arm}} & \textcolor{ChalmersPurple}{\textbf{Attacker processes}} & \textcolor{ChalmersPurple}{\textbf{Attacker writes}} & \textcolor{ChalmersPurple}{\textbf{Honest visible}} & \textcolor{ChalmersPurple}{\textbf{Attacker / list}} & \textcolor{ChalmersPurple}{\textbf{Attacker share}} \\[1mm]\hline
      32 & honest & 0 & 0 & 32/32 & 0/32 & 0.000000 \\
      32 & flood & 1 & 8 & 32/32 & 8/40 & 0.200000 \\
      32 & flood & 1 & 32 & 32/32 & 32/64 & 0.500000 \\
      32 & flood & 1 & 128 & 32/32 & 128/160 & 0.800000 \\
      32 & eclipse & 5 & 1 & 0/32 & 1/1 & 1.000000 \\
    \end{tabular}\endgroup
  };
  \node[quiet,font=\scriptsize,text width=13.3cm] at (7.2,0.65) {Flooding adds advertised identities; custody can refuse honest writes. Neither row measures group capture.};
  \node[font=\scriptsize] at (5.6,0.15) {\hyperlink{table-index}{\beamerbutton{All tables}}};
  \node[font=\scriptsize] at (8.6,0.15) {\hyperlink{table-admission}{\beamerbutton{Next table}}};
\end{visual}
\resultlink{result-custody}{Back to result}
\note{All 15 saved rows; one run per configuration. Cold-key discovery, not group capture. Flooding uses one process publishing multiple advertised identities. Eclipse uses five crafted record holders refusing honest advertisements. Source: research/hivemind/results/step3\_candidates\_live.csv. Exact saved counts; displayed decimals are rounded. Source paths and full per-run calculations: EXPERIMENT\_TABLES.md and experiment-tables-data.json; regenerated by export\_experiment\_tables.py.}
\end{frame}

\begin{frame}[label=table-admission]{Admission: measured group composition}
\begin{visual}
  \node[quiet,text=muted,text width=13.5cm,font=\footnotesize] at (7.2,5.3) {16 training peers; target group 8; minimum 2; one run per row};
  \node[anchor=north west,inner sep=0pt,font=\scriptsize] at (0.45,4.66) {%
    \begingroup\fontsize{8}{10}\selectfont\setlength{\tabcolsep}{3pt}\renewcommand{\arraystretch}{1.0}%
    \setbox\strutbox=\hbox{\vrule height7pt depth3pt width0pt}%
    \begin{tabular}{>{\raggedright\arraybackslash}p{2.7cm}>{\raggedright\arraybackslash}p{0.55cm}>{\raggedright\arraybackslash}p{1.55cm}>{\raggedright\arraybackslash}p{2.35cm}>{\raggedright\arraybackslash}p{2.35cm}>{\raggedright\arraybackslash}p{2.3cm}}
      \textcolor{ChalmersPurple}{\textbf{Arm}} & \textcolor{ChalmersPurple}{\textbf{k}} & \textcolor{ChalmersPurple}{\textbf{Designated share}} & \textcolor{ChalmersPurple}{\textbf{Fully captured}} & \textcolor{ChalmersPurple}{\textbf{Victim + all k}} & \textcolor{ChalmersPurple}{\textbf{Mean beta}} \\[1mm]\hline
      Honest control & 1 & 1/16 & 0/20 & 0/20 & 0.000000 \\
      Selective admission & 1 & 1/16 & 12/21 & 12/21 & 0.315476 \\
      Honest control & 2 & 2/16 & 0/20 & 0/20 & 0.068750 \\
      Selective admission & 2 & 2/16 & 17/20 & 10/20 & 0.508333 \\
      Honest control & 4 & 4/16 & 0/20 & 0/20 & 0.275000 \\
      Selective admission & 4 & 4/16 & 14/21 & 7/21 & 0.554932 \\
    \end{tabular}\endgroup
  };
  \node[quiet,font=\scriptsize,text width=13.3cm] at (7.2,0.65) {Control: k designated peers behave honestly. Attack: full capture = victim + attackers only. Leadership was not logged.};
  \node[font=\scriptsize] at (5.6,0.15) {\hyperlink{table-index}{\beamerbutton{All tables}}};
  \node[font=\scriptsize] at (8.6,0.15) {\hyperlink{table-lookup_threshold}{\beamerbutton{Next table}}};
\end{visual}
\resultlink{result-admission}{Back to result}
\note{16 training peers; target group size 8; minimum 2; one run per row. Full capture = victim + at least one attacker, no other honest peer. In honest controls, the k designated comparison identities behave honestly: shares and beta count those identities, not active adversaries. The 'victim + all k' groups have beta = k/(k+1): 0.500, 0.667, 0.800. Other captured groups can contain fewer attackers. These CSVs do not log the leader; composition is measured, leadership count is not. HANDOFF's 'led captures' column is the victim + all k column, not all full captures. Mean beta is over every recorded victim round. For k=1 it is 5.047619 times alpha, while a captured pair is 8 times alpha. Exact saved counts; displayed decimals are rounded. Source paths and full per-run calculations: EXPERIMENT\_TABLES.md and experiment-tables-data.json; regenerated by export\_experiment\_tables.py.}
\end{frame}

\begin{frame}[label=table-lookup_threshold]{Lookup coverage: exact captured rounds}
\begin{visual}
  \node[quiet,text=muted,text width=13.5cm,font=\footnotesize] at (7.2,5.3) {8 trainers, 3 attackers; lookup width 20; caching on; 15 epochs per run};
  \node[anchor=north west,inner sep=0pt,font=\footnotesize] at (0.45,4.66) {%
    \begingroup\fontsize{8}{10}\selectfont\setlength{\tabcolsep}{3pt}\renewcommand{\arraystretch}{1.0}%
    \setbox\strutbox=\hbox{\vrule height7pt depth3pt width0pt}%
    \begin{tabular}{>{\raggedright\arraybackslash}p{1.8cm}>{\raggedright\arraybackslash}p{4.8cm}>{\raggedright\arraybackslash}p{2.55cm}>{\raggedright\arraybackslash}p{2.6cm}}
      \textcolor{ChalmersPurple}{\textbf{DHT Sybils}} & \textcolor{ChalmersPurple}{\textbf{Captured / rounds, per repeat}} & \textcolor{ChalmersPurple}{\textbf{Mean capture}} & \textcolor{ChalmersPurple}{\textbf{Mean beta}} \\[1mm]\hline
      5 & 0/15, 0/15 & 0.000000\% & 0.350000 \\
      10 & 0/15, 1/16 & 3.125000\% & 0.391667 \\
      15 & 4/15, 0/16 & 13.333333\% & 0.305035 \\
      18 & 3/14, 1/16 & 13.839286\% & 0.436384 \\
      20 & 11/11, 9/9 & 100.000000\% & 0.656566 \\
      22 & 11/11, 14/14 & 100.000000\% & 0.688312 \\
      25 & 15/15, 13/13, 12/12 & 100.000000\% & 0.663711 \\
      40 & 12/12, 10/10 & 100.000000\% & 0.695833 \\
    \end{tabular}\endgroup
  };
  \node[quiet,font=\scriptsize,text width=13.3cm] at (7.2,0.65) {Means weight repeats equally. Sybils joined before training; these rates do not measure setup latency.};
  \node[font=\scriptsize] at (5.6,0.15) {\hyperlink{table-index}{\beamerbutton{All tables}}};
  \node[font=\scriptsize] at (8.6,0.15) {\hyperlink{table-lookup_controls}{\beamerbutton{Next table}}};
\end{visual}
\resultlink{result-lookup}{Back to result}
\note{8 training peers, 3 attackers, target group 4, bucket\_size 20, caching on, 15 training epochs per run. Fractions count full captures: victim + attackers only. Means give equal weight to each repeat, not each round. Additional Sybils are DHT-only processes. They start before training; measured rates are not arrival-to-capture latency. The current saved 25-Sybil runs give mean beta 0.663711, differing from HANDOFF's 0.706 summary. Capture remains 40/40. Exact saved counts; displayed decimals are rounded. Source paths and full per-run calculations: EXPERIMENT\_TABLES.md and experiment-tables-data.json; regenerated by export\_experiment\_tables.py.}
\end{frame}

\begin{frame}[label=table-lookup_controls]{Lookup controls: width, cache, fillers}
\begin{visual}
  \node[quiet,text=muted,text width=13.5cm,font=\footnotesize] at (7.2,5.3) {8 training peers; full capture counts; means weight repeats equally};
  \node[anchor=north west,inner sep=0pt,font=\scriptsize] at (0.45,4.66) {%
    \begingroup\fontsize{8}{10}\selectfont\setlength{\tabcolsep}{3pt}\renewcommand{\arraystretch}{1.0}%
    \setbox\strutbox=\hbox{\vrule height7pt depth3pt width0pt}%
    \begin{tabular}{>{\raggedright\arraybackslash}p{1.4cm}>{\raggedright\arraybackslash}p{1.05cm}>{\raggedright\arraybackslash}p{1.05cm}>{\raggedright\arraybackslash}p{1.05cm}>{\raggedright\arraybackslash}p{1.05cm}>{\raggedright\arraybackslash}p{3.9cm}>{\raggedright\arraybackslash}p{2.2cm}}
      \textcolor{ChalmersPurple}{\textbf{Arm}} & \textcolor{ChalmersPurple}{\textbf{Width}} & \textcolor{ChalmersPurple}{\textbf{Sybils}} & \textcolor{ChalmersPurple}{\textbf{Fillers}} & \textcolor{ChalmersPurple}{\textbf{Cache}} & \textcolor{ChalmersPurple}{\textbf{Captured / rounds}} & \textcolor{ChalmersPurple}{\textbf{Mean capture}} \\[1mm]\hline
      Chase & 20 & 5 & 20 & on & 0/15, 0/16 & 0.000000\% \\
      Chase & 20 & 25 & 0 & off & 7/7, 13/13 & 100.000000\% \\
      Chase & 40 & 25 & 0 & on & 2/16, 1/16 & 9.375000\% \\
      Chase & 40 & 45 & 0 & on & 13/13, 10/10 & 100.000000\% \\
      Honest & 20 & 0 & 0 & on & 0/15, 0/15, 0/16 & 0.000000\% \\
      Honest & 40 & 0 & 0 & on & 0/16 & 0.000000\% \\
      Honest & 20 & 0 & 20 & on & 0/16 & 0.000000\% \\
    \end{tabular}\endgroup
  };
  \node[quiet,font=\scriptsize,text width=13.3cm] at (7.2,0.65) {Idle fillers control process count; they do not establish a larger active routing overlay.};
  \node[font=\scriptsize] at (5.6,0.15) {\hyperlink{table-index}{\beamerbutton{All tables}}};
  \node[font=\scriptsize] at (8.6,0.15) {\hyperlink{table-lookup_history}{\beamerbutton{Next table}}};
\end{visual}
\resultlink{result-lookup}{Back to result}
\note{8 training peers throughout. Honest arms have no malicious trainers; k1 in their filenames is an unused configured attacker count. 20 idle fillers do not establish a larger active routing overlay. Read this as a process-count control. Equal-weight mean of repeat rates. The captured definition is the same as the threshold table. Exact saved counts; displayed decimals are rounded. Source paths and full per-run calculations: EXPERIMENT\_TABLES.md and experiment-tables-data.json; regenerated by export\_experiment\_tables.py.}
\end{frame}

\begin{frame}[label=table-lookup_history]{Earlier lookup controls: width and fillers}
\begin{visual}
  \node[quiet,text=muted,text width=13.5cm,font=\footnotesize] at (7.2,5.3) {Exploratory controls; 8 training peers, 3 attackers, 5 DHT Sybils; caching on};
  \node[anchor=north west,inner sep=0pt,font=\scriptsize] at (0.45,4.66) {%
    \begingroup\fontsize{8}{10}\selectfont\setlength{\tabcolsep}{3pt}\renewcommand{\arraystretch}{1.0}%
    \setbox\strutbox=\hbox{\vrule height7pt depth3pt width0pt}%
    \begin{tabular}{>{\raggedright\arraybackslash}p{1.1cm}>{\raggedright\arraybackslash}p{1.1cm}>{\raggedright\arraybackslash}p{1.1cm}>{\raggedright\arraybackslash}p{1.8cm}>{\raggedright\arraybackslash}p{4.1cm}>{\raggedright\arraybackslash}p{2.4cm}}
      \textcolor{ChalmersPurple}{\textbf{Width}} & \textcolor{ChalmersPurple}{\textbf{Sybils}} & \textcolor{ChalmersPurple}{\textbf{Fillers}} & \textcolor{ChalmersPurple}{\textbf{DHT nodes}} & \textcolor{ChalmersPurple}{\textbf{Captured / rounds}} & \textcolor{ChalmersPurple}{\textbf{Mean capture}} \\[1mm]\hline
      4 & 5 & 0 & 14 & 7/11, 10/10 & 81.818182\% \\
      8 & 5 & 0 & 14 & 0/15 & 0.000000\% \\
      12 & 5 & 0 & 14 & 0/15 & 0.000000\% \\
      20 & 5 & 12 & 26 & 1/15 & 6.666667\% \\
      20 & 5 & 40 & 54 & 0/15 & 0.000000\% \\
      20 & 5 & 90 & 104 & 1/16 & 6.250000\% \\
    \end{tabular}\endgroup
  };
  \node[quiet,font=\scriptsize,text width=13.3cm] at (7.2,0.65) {DHT nodes = trainers + Sybils + bootstrap + fillers. Idle fillers did not enlarge the victim's routing view.};
  \node[font=\scriptsize] at (5.6,0.15) {\hyperlink{table-index}{\beamerbutton{All tables}}};
  \node[font=\scriptsize] at (8.6,0.15) {\hyperlink{table-sustained}{\beamerbutton{Next table}}};
\end{visual}
\resultlink{result-lookup}{Back to result}
\note{Earlier exploratory controls, separate from the main threshold campaign. Eight training peers, three attackers, five DHT Sybils, caching on. Total here counts DHT nodes: eight trainers' DHT nodes + five Sybils + one bootstrap + fillers; it does not count each trainer's internal Python subprocesses. The 104-node run is 1/16 full captures, not evidence of capture across a larger active routing overlay: idle fillers did not enter the victim's routing table. Counts and equal-repeat means from the saved b4/b8/b12/f12/f40/f90 CSVs; no experiments rerun. Exact saved counts; displayed decimals are rounded. Source paths and full per-run calculations: EXPERIMENT\_TABLES.md and experiment-tables-data.json; regenerated by export\_experiment\_tables.py.}
\end{frame}

\begin{frame}[label=table-sustained]{Sustained capture: all three repeats}
\begin{visual}
  \node[quiet,text=muted,text width=13.5cm,font=\footnotesize] at (7.2,5.3) {8 trainers, 3 attackers, 25 DHT Sybils; caching on; sharding off};
  \node[anchor=north west,inner sep=0pt,font=\footnotesize] at (0.45,4.66) {%
    \begingroup\fontsize{8}{10}\selectfont\setlength{\tabcolsep}{3pt}\renewcommand{\arraystretch}{1.0}%
    \setbox\strutbox=\hbox{\vrule height7pt depth3pt width0pt}%
    \begin{tabular}{>{\raggedright\arraybackslash}p{1.25cm}>{\raggedright\arraybackslash}p{4.35cm}>{\raggedright\arraybackslash}p{2.45cm}>{\raggedright\arraybackslash}p{3.6cm}}
      \textcolor{ChalmersPurple}{\textbf{Repeat}} & \textcolor{ChalmersPurple}{\textbf{Captured / recorded rounds}} & \textcolor{ChalmersPurple}{\textbf{Mean beta}} & \textcolor{ChalmersPurple}{\textbf{Group size: rounds}} \\[1mm]\hline
      1 & 15/15 & 0.700000 & 2: 3, 4: 12 \\
      2 & 13/13 & 0.589744 & 2: 7, 3: 4, 4: 2 \\
      3 & 12/12 & 0.701389 & 2: 2, 3: 1, 4: 9 \\
    \end{tabular}\endgroup
  };
  \node[quiet,font=\scriptsize,text width=13.3cm] at (7.2,0.65) {40/40 consecutive recorded victim rounds captured. Fifteen training epochs do not imply fifteen rounds.};
  \node[font=\scriptsize] at (5.6,0.15) {\hyperlink{table-index}{\beamerbutton{All tables}}};
  \node[font=\scriptsize] at (8.6,0.15) {\hyperlink{table-tracking}{\beamerbutton{Next table}}};
\end{visual}
\resultlink{result-sustained}{Back to result}
\note{40/40 recorded victim rounds captured: 15+13+12. The runs are 15 epochs, not necessarily 15 recorded rounds. 8 training peers; 3 attackers; 25 additional DHT Sybils; target group 4; caching on; sharding off. Prepositioned chase policy. This is full capture; it differs from sharing a group in the tracking experiment. Exact saved counts; displayed decimals are rounded. Source paths and full per-run calculations: EXPERIMENT\_TABLES.md and experiment-tables-data.json; regenerated by export\_experiment\_tables.py.}
\end{frame}

\begin{frame}[label=table-tracking]{Tracking: co-membership and prediction}
\begin{visual}
  \node[quiet,text=muted,text width=13.5cm,font=\footnotesize] at (7.2,5.3) {16 peers | 4 shards; target group 4; 4 attackers in attack arms; 2 repeats; 12 epochs};
  \node[anchor=north west,inner sep=0pt,font=\scriptsize] at (0.45,4.66) {%
    \begingroup\fontsize{8}{10}\selectfont\setlength{\tabcolsep}{3pt}\renewcommand{\arraystretch}{1.0}%
    \setbox\strutbox=\hbox{\vrule height7pt depth3pt width0pt}%
    \begin{tabular}{>{\raggedright\arraybackslash}p{0.95cm}>{\raggedright\arraybackslash}p{1.9cm}>{\raggedright\arraybackslash}p{3.6cm}>{\raggedright\arraybackslash}p{2.35cm}>{\raggedright\arraybackslash}p{3.0cm}}
      \textcolor{ChalmersPurple}{\textbf{Peers}} & \textcolor{ChalmersPurple}{\textbf{Arm}} & \textcolor{ChalmersPurple}{\textbf{Shared / rounds, per repeat}} & \textcolor{ChalmersPurple}{\textbf{Mean shared}} & \textcolor{ChalmersPurple}{\textbf{Correct predictions}} \\[1mm]\hline
      16 & none & 0/12, 0/12 & 0.000000\% & -- \\
      16 & naive & 7/12, 7/13 & 56.089744\% & -- \\
      16 & tracking & 9/12, 12/13 & 83.653846\% & 12/12, 13/13 \\
    \end{tabular}\endgroup
  };
  \node[quiet,font=\scriptsize,text width=13.3cm] at (7.2,0.65) {Shared group is not full capture. Predictions require current group information.};
  \node[font=\scriptsize] at (5.6,0.15) {\hyperlink{table-index}{\beamerbutton{All tables}}};
  \node[font=\scriptsize] at (8.6,0.15) {\hyperlink{table-tracking-2}{\beamerbutton{Next table}}};
\end{visual}
\resultlink{result-tracking}{Back to result}
\note{Target group size 4; 4 shards; 4 attackers in naive/tracking arms, none in none arm; 12 epochs; sharding explicitly enabled. Sharing means any attacker in the victim's group, not full capture. Mean is unweighted across the two repeats. Predictions are attacker records matched to the victim by group\_id, not by per-peer round number. Multiple attackers may predict the same group. 63/63 prediction records are correct, conditional on group information being available. This does not measure indefinite tracking after losing contact. Exact saved counts; displayed decimals are rounded. Source paths and full per-run calculations: EXPERIMENT\_TABLES.md and experiment-tables-data.json; regenerated by export\_experiment\_tables.py.}
\end{frame}

\begin{frame}[label=table-tracking-2]{Tracking: co-membership and prediction}
\begin{visual}
  \node[quiet,text=muted,text width=13.5cm,font=\footnotesize] at (7.2,5.3) {32 peers | 4 shards; target group 4; 4 attackers in attack arms; 2 repeats; 12 epochs};
  \node[anchor=north west,inner sep=0pt,font=\scriptsize] at (0.45,4.66) {%
    \begingroup\fontsize{8}{10}\selectfont\setlength{\tabcolsep}{3pt}\renewcommand{\arraystretch}{1.0}%
    \setbox\strutbox=\hbox{\vrule height7pt depth3pt width0pt}%
    \begin{tabular}{>{\raggedright\arraybackslash}p{0.95cm}>{\raggedright\arraybackslash}p{1.9cm}>{\raggedright\arraybackslash}p{3.6cm}>{\raggedright\arraybackslash}p{2.35cm}>{\raggedright\arraybackslash}p{3.0cm}}
      \textcolor{ChalmersPurple}{\textbf{Peers}} & \textcolor{ChalmersPurple}{\textbf{Arm}} & \textcolor{ChalmersPurple}{\textbf{Shared / rounds, per repeat}} & \textcolor{ChalmersPurple}{\textbf{Mean shared}} & \textcolor{ChalmersPurple}{\textbf{Correct predictions}} \\[1mm]\hline
      32 & none & 0/12, 0/12 & 0.000000\% & -- \\
      32 & naive & 5/12, 5/13 & 40.064103\% & -- \\
      32 & tracking & 9/12, 9/13 & 72.115385\% & 14/14, 11/11 \\
    \end{tabular}\endgroup
  };
  \node[quiet,font=\scriptsize,text width=13.3cm] at (7.2,0.65) {Shared group is not full capture. Predictions require current group information.};
  \node[font=\scriptsize] at (5.6,0.15) {\hyperlink{table-index}{\beamerbutton{All tables}}};
  \node[font=\scriptsize] at (8.6,0.15) {\hyperlink{table-tracking-3}{\beamerbutton{Next table}}};
\end{visual}
\resultlink{result-tracking}{Back to result}
\note{Target group size 4; 4 shards; 4 attackers in naive/tracking arms, none in none arm; 12 epochs; sharding explicitly enabled. Sharing means any attacker in the victim's group, not full capture. Mean is unweighted across the two repeats. Predictions are attacker records matched to the victim by group\_id, not by per-peer round number. Multiple attackers may predict the same group. 63/63 prediction records are correct, conditional on group information being available. This does not measure indefinite tracking after losing contact. Exact saved counts; displayed decimals are rounded. Source paths and full per-run calculations: EXPERIMENT\_TABLES.md and experiment-tables-data.json; regenerated by export\_experiment\_tables.py.}
\end{frame}

\begin{frame}[label=table-tracking-3]{Tracking: co-membership and prediction}
\begin{visual}
  \node[quiet,text=muted,text width=13.5cm,font=\footnotesize] at (7.2,5.3) {48 peers | 4 shards; target group 4; 4 attackers in attack arms; 2 repeats; 12 epochs};
  \node[anchor=north west,inner sep=0pt,font=\scriptsize] at (0.45,4.66) {%
    \begingroup\fontsize{8}{10}\selectfont\setlength{\tabcolsep}{3pt}\renewcommand{\arraystretch}{1.0}%
    \setbox\strutbox=\hbox{\vrule height7pt depth3pt width0pt}%
    \begin{tabular}{>{\raggedright\arraybackslash}p{0.95cm}>{\raggedright\arraybackslash}p{1.9cm}>{\raggedright\arraybackslash}p{3.6cm}>{\raggedright\arraybackslash}p{2.35cm}>{\raggedright\arraybackslash}p{3.0cm}}
      \textcolor{ChalmersPurple}{\textbf{Peers}} & \textcolor{ChalmersPurple}{\textbf{Arm}} & \textcolor{ChalmersPurple}{\textbf{Shared / rounds, per repeat}} & \textcolor{ChalmersPurple}{\textbf{Mean shared}} & \textcolor{ChalmersPurple}{\textbf{Correct predictions}} \\[1mm]\hline
      48 & none & 0/12, 0/12 & 0.000000\% & -- \\
      48 & naive & 2/13, 3/14 & 18.406593\% & -- \\
      48 & tracking & 3/13, 7/13 & 38.461538\% & 3/3, 10/10 \\
    \end{tabular}\endgroup
  };
  \node[quiet,font=\scriptsize,text width=13.3cm] at (7.2,0.65) {Shared group is not full capture. Predictions require current group information.};
  \node[font=\scriptsize] at (5.6,0.15) {\hyperlink{table-index}{\beamerbutton{All tables}}};
  \node[font=\scriptsize] at (8.6,0.15) {\hyperlink{table-poisoning}{\beamerbutton{Next table}}};
\end{visual}
\resultlink{result-tracking}{Back to result}
\note{Target group size 4; 4 shards; 4 attackers in naive/tracking arms, none in none arm; 12 epochs; sharding explicitly enabled. Sharing means any attacker in the victim's group, not full capture. Mean is unweighted across the two repeats. Predictions are attacker records matched to the victim by group\_id, not by per-peer round number. Multiple attackers may predict the same group. 63/63 prediction records are correct, conditional on group information being available. This does not measure indefinite tracking after losing contact. Exact saved counts; displayed decimals are rounded. Source paths and full per-run calculations: EXPERIMENT\_TABLES.md and experiment-tables-data.json; regenerated by export\_experiment\_tables.py.}
\end{frame}

\begin{frame}[label=table-poisoning]{Poisoning: distance to the optimum}
\begin{visual}
  \node[quiet,text=muted,text width=13.5cm,font=\footnotesize] at (7.2,5.3) {Epoch 20; mean Euclidean distance over honest peers; smaller is better};
  \node[anchor=north west,inner sep=0pt,font=\footnotesize] at (0.45,4.66) {%
    \begingroup\fontsize{8}{10}\selectfont\setlength{\tabcolsep}{3pt}\renewcommand{\arraystretch}{1.0}%
    \setbox\strutbox=\hbox{\vrule height7pt depth3pt width0pt}%
    \begin{tabular}{>{\raggedright\arraybackslash}p{5.6cm}>{\raggedright\arraybackslash}p{2.1cm}>{\raggedright\arraybackslash}p{2.1cm}>{\raggedright\arraybackslash}p{2.1cm}}
      \textcolor{ChalmersPurple}{\textbf{Configuration}} & \textcolor{ChalmersPurple}{\textbf{Seed 0}} & \textcolor{ChalmersPurple}{\textbf{Seed 1}} & \textcolor{ChalmersPurple}{\textbf{Seed 2}} \\[1mm]\hline
      No attack / mean & 0.672075 & 0.548488 & 0.743582 \\
      Sign flip, weight 1 / mean & 0.540327 & -- & -- \\
      Sign flip, weight $10^6$ / mean & 30.708980 & 81.437714 & 26.034575 \\
      Sign flip, weight $10^6$ / median & 0.619387 & 1.269958 & 0.372256 \\
      Honest gradient, weight $10^6$ / mean & 0.822815 & -- & -- \\
      Honest gradient, weight $10^6$ / median & 0.682287 & -- & -- \\
    \end{tabular}\endgroup
  };
  \node[quiet,font=\scriptsize,text width=13.3cm] at (7.2,0.65) {8 trainers, 1 poisoner, target group 4; actual beta was not logged. -- = unmeasured.};
  \node[font=\scriptsize] at (5.6,0.15) {\hyperlink{table-index}{\beamerbutton{All tables}}};
  \node[font=\scriptsize] at (8.6,0.15) {\hyperlink{table-poisoning_ratios}{\beamerbutton{Next table}}};
\end{visual}
\resultlink{result-poisoning}{Back to result}
\note{Epoch-20 Euclidean distance to the known optimum, averaged over honest peers (8 control, 7 with attacker). Smaller is better. 8 training peers; one poisoning peer; target group size 4. Actual per-round group composition was not logged or forced; beta=1/4 is not a measured constant. 64-parameter linear regression; sign flip sends -5 times the honest gradient. Missing cells were not measured. Ratios use each seed's own no-attack control. The separate step2 0.1492 baseline is not comparable. Exact saved counts; displayed decimals are rounded. Source paths and full per-run calculations: EXPERIMENT\_TABLES.md and experiment-tables-data.json; regenerated by export\_experiment\_tables.py.}
\end{frame}

\begin{frame}[label=table-poisoning_ratios]{Poisoning: ratios to matched controls}
\begin{visual}
  \node[quiet,text=muted,text width=13.5cm,font=\footnotesize] at (7.2,5.3) {Honest peers' epoch-20 distance / no-attack distance from the same seed};
  \node[anchor=north west,inner sep=0pt,font=\footnotesize] at (0.45,4.66) {%
    \begingroup\fontsize{8}{10}\selectfont\setlength{\tabcolsep}{3pt}\renewcommand{\arraystretch}{1.0}%
    \setbox\strutbox=\hbox{\vrule height7pt depth3pt width0pt}%
    \begin{tabular}{>{\raggedright\arraybackslash}p{5.6cm}>{\raggedright\arraybackslash}p{2.1cm}>{\raggedright\arraybackslash}p{2.1cm}>{\raggedright\arraybackslash}p{2.1cm}}
      \textcolor{ChalmersPurple}{\textbf{Configuration}} & \textcolor{ChalmersPurple}{\textbf{Seed 0}} & \textcolor{ChalmersPurple}{\textbf{Seed 1}} & \textcolor{ChalmersPurple}{\textbf{Seed 2}} \\[1mm]\hline
      No attack / mean & 1.000000 & 1.000000 & 1.000000 \\
      Sign flip, weight 1 / mean & 0.803967 & -- & -- \\
      Sign flip, weight $10^6$ / mean & 45.692765 & 148.476664 & 35.012399 \\
      Sign flip, weight $10^6$ / median & 0.921604 & 2.315378 & 0.500625 \\
      Honest gradient, weight $10^6$ / mean & 1.224290 & -- & -- \\
      Honest gradient, weight $10^6$ / median & 1.015194 & -- & -- \\
    \end{tabular}\endgroup
  };
  \node[quiet,font=\scriptsize,text width=13.3cm] at (7.2,0.65) {1 = matched control. Uses step7 controls, not the separate step2 baseline. -- = unmeasured.};
  \node[font=\scriptsize] at (5.6,0.15) {\hyperlink{table-index}{\beamerbutton{All tables}}};
  \node[font=\scriptsize] at (8.6,0.15) {\hyperlink{table-withdrawal}{\beamerbutton{Next table}}};
\end{visual}
\resultlink{result-poisoning}{Back to result}
\note{Epoch-20 Euclidean distance to the known optimum, averaged over honest peers (8 control, 7 with attacker). Smaller is better. 8 training peers; one poisoning peer; target group size 4. Actual per-round group composition was not logged or forced; beta=1/4 is not a measured constant. 64-parameter linear regression; sign flip sends -5 times the honest gradient. Missing cells were not measured. Ratios use each seed's own no-attack control. The separate step2 0.1492 baseline is not comparable. Exact saved counts; displayed decimals are rounded. Source paths and full per-run calculations: EXPERIMENT\_TABLES.md and experiment-tables-data.json; regenerated by export\_experiment\_tables.py.}
\end{frame}

\begin{frame}[label=table-withdrawal]{Withdrawal: no clear capture gain}
\begin{visual}
  \node[quiet,text=muted,text width=13.5cm,font=\footnotesize] at (7.2,5.3) {8 training peers; 1 attacker; target group 4; caching on; 3 repeats per arm};
  \node[anchor=north west,inner sep=0pt,font=\footnotesize] at (0.45,4.66) {%
    \begingroup\fontsize{8}{10}\selectfont\setlength{\tabcolsep}{3pt}\renewcommand{\arraystretch}{1.0}%
    \setbox\strutbox=\hbox{\vrule height7pt depth3pt width0pt}%
    \begin{tabular}{>{\raggedright\arraybackslash}p{3.4cm}>{\raggedright\arraybackslash}p{5.1cm}>{\raggedright\arraybackslash}p{3.3cm}}
      \textcolor{ChalmersPurple}{\textbf{Arm}} & \textcolor{ChalmersPurple}{\textbf{Captured / rounds, per repeat}} & \textcolor{ChalmersPurple}{\textbf{Mean capture}} \\[1mm]\hline
      Leader only & 11/16, 5/16, 12/16 & 58.333333\% \\
      Leader + withdrawal & 7/16, 14/16, 9/16 & 62.500000\% \\
    \end{tabular}\endgroup
  };
  \node[quiet,font=\scriptsize,text width=13.3cm] at (7.2,0.65) {Within-arm variation exceeds the between-arm gap. No clear additional capture benefit.};
  \node[font=\scriptsize] at (5.6,0.15) {\hyperlink{table-index}{\beamerbutton{All tables}}};
  \node[font=\scriptsize] at (8.6,0.15) {\hyperlink{table-withdrawal_discovery}{\beamerbutton{Next table}}};
\end{visual}
\resultlink{result-limitations}{Back to result}
\note{8 training peers; 1 attacker; target group 4; caching on; 3 repeats per arm. Full capture definition as above; the between-arm gap is much smaller than variation across repeats. This is not a successful additional group-capture mechanism; the separate DHT censorship effect is narrower. Exact saved counts; displayed decimals are rounded. Source paths and full per-run calculations: EXPERIMENT\_TABLES.md and experiment-tables-data.json; regenerated by export\_experiment\_tables.py.}
\end{frame}

\begin{frame}[label=table-withdrawal_discovery]{Withdrawal: discovery suppression}
\begin{visual}
  \node[quiet,text=muted,text width=13.5cm,font=\footnotesize] at (7.2,5.3) {DHT-only experiment; two repeats per setting; no training};
  \node[anchor=north west,inner sep=0pt,font=\scriptsize] at (0.45,4.66) {%
    \begingroup\fontsize{8}{10}\selectfont\setlength{\tabcolsep}{3pt}\renewcommand{\arraystretch}{1.0}%
    \setbox\strutbox=\hbox{\vrule height7pt depth3pt width0pt}%
    \begin{tabular}{>{\raggedright\arraybackslash}p{1.65cm}>{\raggedright\arraybackslash}p{1.45cm}>{\raggedright\arraybackslash}p{1.3cm}>{\raggedright\arraybackslash}p{2.3cm}>{\raggedright\arraybackslash}p{2.6cm}>{\raggedright\arraybackslash}p{2.5cm}}
      \textcolor{ChalmersPurple}{\textbf{Honest peers}} & \textcolor{ChalmersPurple}{\textbf{Targeted}} & \textcolor{ChalmersPurple}{\textbf{Repeats}} & \textcolor{ChalmersPurple}{\textbf{Hidden / repeat}} & \textcolor{ChalmersPurple}{\textbf{Visible survivors}} & \textcolor{ChalmersPurple}{\textbf{Accepted writes}} \\[1mm]\hline
      12 & 1 & 2 & 1, 1 & 11, 11 & 1, 1 \\
      12 & 3 & 2 & 3, 3 & 9, 9 & 3, 3 \\
      12 & 6 & 2 & 6, 6 & 6, 6 & 6, 6 \\
      12 & 11 & 2 & 11, 11 & 1, 1 & 11, 11 \\
      40 & 1 & 2 & 1, 1 & 39, 39 & 1, 1 \\
      40 & 20 & 2 & 20, 20 & 20, 20 & 20, 20 \\
      40 & 39 & 2 & 39, 39 & 1, 1 & 39, 39 \\
    \end{tabular}\endgroup
  };
  \node[quiet,font=\scriptsize,text width=13.3cm] at (7.2,0.65) {Suppressing an advertisement did not establish an additional group-capture mechanism.};
  \node[font=\scriptsize] at (5.6,0.15) {\hyperlink{table-index}{\beamerbutton{All tables}}};
  \node[font=\scriptsize] at (8.6,0.15) {\hyperlink{table-baseline}{\beamerbutton{Next table}}};
\end{visual}
\resultlink{result-limitations}{Back to result}
\note{Source: research/hivemind/results/step8\_withdrawal\_phase2.csv and step8\_withdrawal\_phase2\_n40.csv. No training in these runs. Every intended advertisement was suppressed, every non-target remained. The paired training experiment above finds no clear added capture gain. At both sizes, separate fightback rows record zero successful visibility restorations across 30 redeclarations / 30 s. Exact saved counts; displayed decimals are rounded. Source paths and full per-run calculations: EXPERIMENT\_TABLES.md and experiment-tables-data.json; regenerated by export\_experiment\_tables.py.}
\end{frame}

\begin{frame}[label=table-baseline]{Baseline training: recorded endpoints}
\begin{visual}
  \node[quiet,text=muted,text width=13.5cm,font=\footnotesize] at (7.2,5.3) {64-parameter regression; one run per row; \_b = independent repeat};
  \node[anchor=north west,inner sep=0pt,font=\scriptsize] at (0.45,4.66) {%
    \begingroup\fontsize{8}{10}\selectfont\setlength{\tabcolsep}{3pt}\renewcommand{\arraystretch}{1.0}%
    \setbox\strutbox=\hbox{\vrule height7pt depth3pt width0pt}%
    \begin{tabular}{>{\raggedright\arraybackslash}p{2.0cm}>{\raggedright\arraybackslash}p{1.05cm}>{\raggedright\arraybackslash}p{1.8cm}>{\raggedright\arraybackslash}p{2.4cm}>{\raggedright\arraybackslash}p{2.4cm}>{\raggedright\arraybackslash}p{2.25cm}}
      \textcolor{ChalmersPurple}{\textbf{Peers x group}} & \textcolor{ChalmersPurple}{\textbf{Epoch}} & \textcolor{ChalmersPurple}{\textbf{Peers recorded}} & \textcolor{ChalmersPurple}{\textbf{Mean distance}} & \textcolor{ChalmersPurple}{\textbf{Peer SD}} & \textcolor{ChalmersPurple}{\textbf{Min. peers seen}} \\[1mm]\hline
      16x4 & 20 & 16 & 0.126592 & 0.000595 & 16 \\
      16x8 & 20 & 16 & 0.122636 & 0.000725 & 16 \\
      16x8\_b & 20 & 16 & 0.122837 & 0.000345 & 16 \\
      24x12 & 20 & 24 & 0.122883 & 0.000923 & 24 \\
      24x8 & 20 & 24 & 0.123333 & 0.001474 & 24 \\
      24x8\_b & 20 & 24 & 0.123384 & 0.001240 & 24 \\
      8x8 & 20 & 8 & 0.183949 & 0.000000 & 8 \\
    \end{tabular}\endgroup
  };
  \node[quiet,font=\scriptsize,text width=13.3cm] at (7.2,0.65) {Different group-size configuration: these are not matched controls for the poisoning runs.};
  \node[font=\scriptsize] at (5.6,0.15) {\hyperlink{table-index}{\beamerbutton{All tables}}};
  \node[font=\scriptsize] at (8.6,0.15) {\hyperlink{table-mixing}{\beamerbutton{Next table}}};
\end{visual}
\resultlink{result-limitations}{Back to result}
\note{64-parameter linear regression; one run per row; \_b is an independent repeat. step2's gradient target group size was unbound. These endpoints are not matched controls for step3/6/7 scripts. The current saved 8x8 endpoint is 0.183949, differing from HANDOFF's 0.1492 canonical-run summary; use the CSV value here. Source files: step2\_training\_*.csv. Exact saved counts; displayed decimals are rounded. Source paths and full per-run calculations: EXPERIMENT\_TABLES.md and experiment-tables-data.json; regenerated by export\_experiment\_tables.py.}
\end{frame}

\begin{frame}[label=table-mixing]{Baseline membership: repeated partners}
\begin{visual}
  \node[quiet,text=muted,text width=13.5cm,font=\footnotesize] at (7.2,5.3) {Observed membership statistics; sharding off; no Monte Carlo rerun};
  \node[anchor=north west,inner sep=0pt,font=\footnotesize] at (0.45,4.66) {%
    \begingroup\fontsize{8}{10}\selectfont\setlength{\tabcolsep}{3pt}\renewcommand{\arraystretch}{1.0}%
    \setbox\strutbox=\hbox{\vrule height7pt depth3pt width0pt}%
    \begin{tabular}{>{\raggedright\arraybackslash}p{2.1cm}>{\raggedright\arraybackslash}p{1.2cm}>{\raggedright\arraybackslash}p{3.0cm}>{\raggedright\arraybackslash}p{2.65cm}>{\raggedright\arraybackslash}p{2.65cm}}
      \textcolor{ChalmersPurple}{\textbf{Peers x group}} & \textcolor{ChalmersPurple}{\textbf{Rounds}} & \textcolor{ChalmersPurple}{\textbf{Pair co-occurrence SD}} & \textcolor{ChalmersPurple}{\textbf{Most shared rounds}} & \textcolor{ChalmersPurple}{\textbf{Distinct partitions}} \\[1mm]\hline
      16x4 & 20 & 4.959839 & 17 & 15/20 \\
      16x8 & 20 & 5.527708 & 20 & 13/20 \\
      16x8\_b & 20 & 5.451809 & 20 & 10/20 \\
      24x12 & 20 & 2.239166 & 15 & 20/20 \\
      24x8 & 20 & 1.961507 & 12 & 20/20 \\
      24x8\_b & 20 & 2.008959 & 14 & 20/20 \\
    \end{tabular}\endgroup
  };
  \node[quiet,font=\scriptsize,text width=13.3cm] at (7.2,0.65) {Generalisation is unsettled. Invalid 8x8 partition alignment is excluded; see notes.};
  \node[font=\scriptsize] at (5.6,0.15) {\hyperlink{table-index}{\beamerbutton{All tables}}};
  \node[font=\scriptsize] at (8.6,0.15) {\hyperlink{table-sharding_handoff}{\beamerbutton{Next table}}};
\end{visual}
\resultlink{result-limitations}{Back to result}
\note{Observed statistics recomputed directly from saved group memberships, using the definitions in analyse\_groups.py. No Monte Carlo rerun. HANDOFF 4.8 reports rounded null p-values: 16-peer runs <5e-5; 24x8 0.85/0.70; 24x12 0.46. Exact simulation output was not found among saved CSVs. Sharding disabled. Whether the 16-peer clique or 24-peer mixing regime generalises is unsettled. 8x8 omitted: an undersized group left one peer with 19 rounds and others with 20, so per-peer round counters cannot identify a common partition. Exact saved counts; displayed decimals are rounded. Source paths and full per-run calculations: EXPERIMENT\_TABLES.md and experiment-tables-data.json; regenerated by export\_experiment\_tables.py.}
\end{frame}

\begin{frame}[label=table-sharding_handoff]{Sharding failures: reported observations}
\begin{visual}
  \node[quiet,text=ctl,text width=13.5cm,font=\footnotesize] at (7.2,5.3) {Reported only; raw output unavailable};
  \node[anchor=north west,inner sep=0pt,font=\footnotesize] at (0.45,4.66) {%
    \begingroup\fontsize{8}{10}\selectfont\setlength{\tabcolsep}{3pt}\renewcommand{\arraystretch}{1.0}%
    \setbox\strutbox=\hbox{\vrule height7pt depth3pt width0pt}%
    \begin{tabular}{>{\raggedright\arraybackslash}p{5.2cm}>{\raggedright\arraybackslash}p{7.0cm}}
      \textcolor{ChalmersPurple}{\textbf{Reported configuration}} & \textcolor{ChalmersPurple}{\textbf{Reported observation}} \\[1mm]\hline
      16 peers, group 4 & 8/16 never joined a group \\
      16 peers, group 2 & 12/16 never joined a group \\
      Matchmaking window 15 s & 16/16 participated; zero failures; 96 s vs 337 s \\
      Longer sharded run & Clean at 4 epochs; 63 failures by epoch 10 \\
    \end{tabular}\endgroup
  };
  \node[quiet,font=\scriptsize,text width=13.3cm] at (7.2,0.65) {HANDOFF 4.9; values could not be independently reconstructed from saved run tables.};
  \node[font=\scriptsize] at (5.6,0.15) {\hyperlink{table-index}{\beamerbutton{All tables}}};
  \node[font=\scriptsize] at (8.6,0.15) {\hyperlink{table-index}{\beamerbutton{Next table}}};
\end{visual}
\resultlink{result-limitations}{Back to result}
\note{Source is HANDOFF 4.9 only. Matching per-run raw output tables were not located; these values were not independently recomputed. These are reported exploratory results, not additional validated CSV measurements. Exact saved counts; displayed decimals are rounded. Source paths and full per-run calculations: EXPERIMENT\_TABLES.md and experiment-tables-data.json; regenerated by export\_experiment\_tables.py.}
\end{frame}

% \section{Backup}
% % Backup: admission and aggregation guarantees
\begin{frame}[label=backup-robustness]{Admission changes the group size}
\begin{visual}
  \node[heading] at (3.05,5.05) {Nominal defence setting};
  \node[heading] at (11.1,5.05) {A captured pair};
  \draw[ChalmersPurple,dashed,rounded corners=5mm] (0.7,1.9) rectangle (5.4,4.25);
  \node[victim] at (1.65,3.55) {V};
  \node[peer] at (3.1,3.55) {};
  \node[peer] at (4.55,3.55) {};
  \node[attacker] at (3.1,2.5) {A};
  \node[labeltext,text=ChalmersPurple] at (3.05,1.2) {3 honest : 1 attacker};
  \draw[control] (5.75,3.1)--(8.4,3.1);
  \node[quiet] at (7.1,4.0) {selective admission};
  \draw[ctl,dashed,rounded corners=5mm] (8.9,1.9) rectangle (13.35,4.25);
  \node[victim,minimum size=12mm] at (10.0,3.1) {V};
  \node[attacker,minimum size=14mm] at (12.2,3.1) {A};
  \node[labeltext,text=ctl] at (11.1,1.2) {1 honest : 1 attacker};
  \node[labeltext,font=\normalsize\bfseries] at (7.2,0.2) {The honest-majority condition is gone.};
\end{visual}
\resultlink{result-admission}{Back to admission}
\note{The left drawing is a nominal four-member group containing the poisoner: the median
experiment targeted this size but did not log every actual group. The separate admission
result can give the victim a two-member group with one attacker, changing numerator and
denominator. Median's strict honest-majority condition is absent. Krum's sufficient condition
is 2f+2<n; n=2,f=1 violates it too. Losing sufficient guarantees does not prove every rule
necessarily fails. The arrow connects separately measured settings, not an end-to-end
measured composition. The research focus is Hivemind's election and admission. Prior work
already studies adversarial peer sampling and local Byzantine density, so do not claim the
broad concern is new. Source: HANDOFF 4.3,4.11,6.
\par Negative result: forged withdrawals did not detectably improve capture: 58 percent
leader-only versus 62 percent with forgery, three repeats per arm, with much larger variation.
Hidden peers still read the key and request admission. Source: HANDOFF 4.12.}
\end{frame}


\end{document}
